%%%%%%%%%%%%%%%%%%%%%%%%%%%%%%%%%%%%%%%%%%%%%%%%%%%%%%%%%%%%%%%%%%%%%%
% Lớp 12 | Phương trình mặt phẳng, đường thẳng, mặt cầu | Phương trình mặt cầu trong không gian
% Dạng 2H5?3-2: Xác định tâm, bán kính, đường kính mặt cầu
% Số câu: 114  (sinh tự động bởi tach_cau.py; từ nay thầy sửa trực tiếp file này)
%%%%%%%%%%%%%%%%%%%%%%%%%%%%%%%%%%%%%%%%%%%%%%%%%%%%%%%%%%%%%%%%%%%%%%

%%%%%%%%%%%%%%%% TRẮC NGHIỆM 4 PHƯƠNG ÁN %%%%%%%%%%%%%%%%

% nguồn: Phan1_Lop_12_Hinh_hoc_Chuong_5
\begin{ex}%[2H5N3-2]%[ma:7eadd4dc]
	Trong không gian $Oxyz$, cho  mặt cầu $(S)\colon(x-1)^2+(y-2)^2+(z-3)^2=16$. Bán kính của mặt cầu bằng
	\choice 
		 {$16$}
		 {$2$}
		 {\True $4$}
		{$8$} 
	\loigiai{Bán kính của mặt  cầu là $R=4$.}
\end{ex}

% nguồn: Phan1_Lop_12_Hinh_hoc_Chuong_5
\begin{ex}%[2H5N3-2]%[ma:15bb5a26]
	Trong không gian $Oxyz$, cho mặt cầu $(S)\colon  (x-3)^2+(y+1)^2+(z-2)^2=4$. Xác định tọa độ tâm $I$ và bán kính $R$ của mặt cầu $(S)$.
	 \choice
		{\True $I(3;-1;2)$, $R=2$} 
		{$I(-3;1;-2)$,  $R=2$}
		{$I(-3;1;-2)$,  $R=4$}
		{$I(3;-1;2)$, $R=4$} 
	\loigiai{ 
		Mặt cầu $(S)$ có tâm  $I(3;-1;2)$ và bán kính $R=2$.
	 }
\end{ex}

% nguồn: Phan1_Lop_12_Hinh_hoc_Chuong_5
\begin{ex}%[2H5N3-2]%[ma:35dc4317]
	Trong  không gian $Oxyz$, cho phương trình mặt cầu $(S)\colon (x-2)^2+(y-2)^2+(z+1)^2=1$. Tọa độ tâm $I$ của mặt cầu $(S)$ là
	 \choice
		 {\True $I(2;2;-1)$}
		 {$I(-2;-2;1)$}
		{$I(-2;2;-1)$} 
		{$I(4;4;-2)$} 
	\loigiai{ 
		Mặt cầu tâm $I(a;b;c)$, bán kính $R$ có phương trình:  $(x-a)^2+(y-b)^2+(z-c)^2=R^2$.\\
		Từ  phương trình mặt cầu đề bài cho $\Rightarrow\heva{&a=2\\&b=2\\&c=-1.}$\\
		Vậy tâm mặt cầu là $I(2;2;-1)$. 
	 }
\end{ex}

% nguồn: Phan1_Lop_12_Hinh_hoc_Chuong_5
\begin{ex}%[2H5N3-2]%[ma:4660b8e0]
	 Trong không gian $Oxyz$, cho mặt cầu $(S)\colon x^2+y^2+z^2-2x-4y+2z-2=0$. Xác định toạ độ tâm $I$ và bán kính $R$ của mặt cầu $(S)$.
	 \choice
		{\True  $I(1;2;-1)$, $R=2 \sqrt{2}$}
		{$I(-1;-2;1)$,  $R=2 \sqrt{2}$}
		 {$I(2;4;-2)$, $R=\sqrt{2}$}
		 {$I(2;4;2)$, $R=8$}
	\loigiai{ 
		 Mặt cầu $(S)$ có tâm $I(1;2;-1)$ và bán kính $R=\sqrt{1^2+2^2+(-1)^2-(-2)}=2 \sqrt{2}$.
	} 
\end{ex}

% nguồn: Phan1_Lop_12_Hinh_hoc_Chuong_5
\begin{ex}%[2H5N3-2]%[ma:851994d7]
	Trong không gian $Oxyz$, cho mặt cầu $(x-1)^2+(y-2)^2+(z+5)^2=16$. Tọa độ  tâm $I$ và bán kính $R$ của mặt cầu~là
	\choice 
		 {\True $I(1;2;-5)$; $R=4$}
		 {$I(1;2;-5)$; $R=16$}
		{$I(-1;-2;5)$; $R=4$} 
		 {$I(-1;-2;5)$; $R=16$}
	 \loigiai{
	Mặt cầu $(x-1)^2+(y-2)^2+(z+5)^2=16$ có tâm và bán kính lần lượt~là $I(1;2;-5)$;  $R=4$.
	} 
\end{ex}

% nguồn: Phan1_Lop_12_Hinh_hoc_Chuong_5
\begin{ex}%[2H5N3-2]%[ma:14c4336e]
	Trong không gian $Oxyz$, cho mặt cầu $(S)\colon x^2+y^2+z^2-2x+4y+6z-6=0$. Tâm của mặt cầu  có tọa độ là
	 \choice
		{\True $(1;-2;-3)$} 
		{$(-2;4;0)$} 
		 {$(-1;2;-3)$}
		{$(-1;2;3)$} 
	\loigiai{ 
		Ta có  $x^2+y^2+z^2-2x+4y+6z-6=0 \Leftrightarrow (x-1)^2+(y+2)^2+(z+3)^2=20$.\\
		Suy ra tâm mặt cầu là $I(1;-2;-3)$. 
	} 
\end{ex}

% nguồn: Phan1_Lop_12_Hinh_hoc_Chuong_5
\begin{ex}%[2H5N3-2]%[ma:9857e9a5]
	Trong không gian $Oxyz$, mặt cầu $(S)$ có phương  trình $x^2+y^2+z^2-\sqrt{2}x+y-z-1=0$. Tính bán kính $R$ của mặt cầu $(S)$.
	 \choice
		{$3$} 
		 {$2$}
		 {\True $R=\sqrt{2}$}
		{$R=\sqrt{3}$} 
	\loigiai{ 
		Tạo độ tâm của mặt cầu $(S)$ là $I\left(\dfrac{\sqrt{2}}{2};-\dfrac{1}{2}; \dfrac{1}{2}\right)$ và hệ số $d=-1$.\\ 
		Bán kính  $R=\sqrt{\left(\dfrac{\sqrt{2}}{2}\right)^2+\left(-\dfrac{1}{2}\right)^2+\left(\dfrac{1}{2}\right)^2-(-1)}=\sqrt{2}>0$.\\
		Vậy  bán kính mặt cầu $(S)$ là $R=\sqrt{2}$.
	} 
\end{ex}

% nguồn: Phan1_Lop_12_Hinh_hoc_Chuong_5
\begin{ex}%[2H5N3-2]%[ma:5609857a]
	 Trong không gian với hệ tọa độ $Oxyz$, mặt cầu $(x-5)^2+(y+8)^2+(z-13)^2=9^2$ có tọa độ tâm là
	\choice 
		{$(5;8;13)$} 
		 {\True $(5;-8;13)$}
		 {$(-5;8;-13)$}
		 {$(-5;-8;-13)$}
	\loigiai{ 
		Mặt cầu $(x-5)^2+(y+8)^2+(z-13)^2=9^2$ có  tọa độ tâm là $(5;-8;13)$.
	} 
\end{ex}

% nguồn: Phan1_Lop_12_Hinh_hoc_Chuong_5
\begin{ex}%[2H5N3-2]%[ma:0ce639cb]
	Trong không gian  $Oxyz$, cho mặt cầu $(S)\colon(x-1)^2+(y-2)^2+(z+3)^2=36$. Tìm tọa độ tâm $I$ và bán kính $R$ của mặt cầu $(S)$.
	\choice 
		{$I(-1;-2;3)$  và $R=6$}
		 {$I(-1;-2;3)$ và $R=36$}
		 {$I(1;2;-3)$ và $R=36$}
		{\True  $I(1;2;-3)$ và $R=6$}
	 \loigiai{
		Mặt cầu $(S)$ có tâm là $I(1;2;-3)$ và bán kính  là $R=\sqrt{36}=6$.
	} 
\end{ex}

% nguồn: Phan1_Lop_12_Hinh_hoc_Chuong_5
\begin{ex}%[2H5N3-2]%[ma:d60f0cfe]
	 Trong không gian $Oxyz$, cho mặt cầu $(S) \colon x^2 +(y+3)^2 +(z-1)^2 =4$. Tọa độ tâm $I$ và bán kính $R$ của $(S)$ lần lượt~là
	\choice 
		{$I(0;-3;1)$, $R=1$} 
		{$I(0;3;-1)$, $R=4$} 
		{\True  $I(0;-3;1)$, $R=2$}
		{$I(0;3;-1)$, $R=2$} 
	\loigiai{ 
		Mặt cầu $(S) \colon x^2 +(y+3)^2 +(z-1)^2 =4$ có tâm $I(0;-3;1)$ và bán  kính $R=2$.
	} 
\end{ex}

% nguồn: Phan1_Lop_12_Hinh_hoc_Chuong_5
\begin{ex}%[2H5N3-2]%[ma:d741b823]
	Trong không gian $Oxyz$, cho mặt cầu  $(S)\colon x^2 + (y+2)^2 + (z-1)^2 = 16$. Tọa độ tâm $I$ của $(S)$ là
	\choice 
		{$I(1;-2;0)$} 
		{$I(1;2;0)$} 
		 {$I(0;2;-1)$}
		 {\True $I(0;-2;1)$}
	\loigiai{ 
		Mặt cầu $(S)\colon  x^2 + (y+2)^2 + (z-1)^2 = 16$ có tâm là $I(0;-2;1)$.
	} 
\end{ex}

% nguồn: Phan1_Lop_12_Hinh_hoc_Chuong_5
\begin{ex}%[2H5N3-2]%[ma:83c8a54c]
	Trong không gian với hệ tọa độ $Oxyz$, cho mặt cầu  $(S) \colon (x+2)^2+(y+3)^2+(z-5)^2=36$ có tọa độ tâm $I$ là
	\choice 
		{$I  \left(1;\dfrac{3}{2};-\dfrac{5}{2}\right)$}
		{$I \left(-1;-\dfrac{3}{2};\dfrac{5}{2}\right)$} 
		 {$I \left(2;3;-5\right)$}
		{\True $I  \left(-2;-3;5\right)$}
	\loigiai{Mặt cầu $(S) \colon (x+2)^2+(y+3)^2+(z-5)^2=36$ có tọa độ tâm $I$ là $I \left(-2;-3;5\right)$.} 
\end{ex}

% nguồn: Phan1_Lop_12_Hinh_hoc_Chuong_5
\begin{ex}%[2H5N3-2]%[ma:4ad4deb8]
	Trong không gian $Oxyz$, cho điểm $A(-1;2;-3)$, $B(3;-4;1)$. Khi  đó mặt cầu $(S)$ nhận $AB$ làm đường kính, có tọa độ tâm $I$ là
	\choice 
		 {$(4;-6;4)$}
		{\True  $(1;-1;-1)$}
		{$(2;-2;-2)$} 
		 {$(2;-3;2)$}
	\loigiai{Gọi $I$ là trung điểm của $AB$. Khi đó $I(1;-1;-1)$ là tâm của  mặt cầu $(S)$.}
\end{ex}

% nguồn: Phan1_Lop_12_Hinh_hoc_Chuong_5
\begin{ex}%[2H5N3-2]%[ma:cc19d726]
	Trong không gian với hệ tọa độ $Oxyz$,  trong các mặt cầu dưới đây, mặt cầu nào có bán kính $R=2$?
	 \choice
		{$(S)\colon x^2+y^2+z^2-4x+2y+2z-3=0$} 
		{$(S)\colon x^2+y^2+z^2-4x+2y+2z-10=0$} 
		{\True $(S)\colon  x^2+y^2+z^2-4x+2y+2z+2=0$}
		{$(S)\colon  x^2+y^2+z^2-4x+2y+2z+5=0$}
	\loigiai{ 
		Ta có mặt cầu $(S)\colon x^2+y^2+z^2-2ax-2by-2cz+d=0$  có bán kính là $R=\sqrt{a^2+b^2+c^2-d}$.\\
		 Trong phương trình $(S)\colon x^2+y^2+z^2-4x+2y+2z+2=0$, ta có $\heva{&a=2\\&b=-1\\&c=-1\\&d=2} \Rightarrow R=\sqrt{a^2+b^2+c^2-d}=\sqrt{4}=2$.
	 }
\end{ex}

% nguồn: Phan1_Lop_12_Hinh_hoc_Chuong_5
\begin{ex}%[2H5N3-2]%[ma:c711589d]
	Trong không gian $Oxyz$, tâm của mặt cầu $(S):(x-2)^2+(y+3)^2+(z-1)^2=16$ có tọa  độ là
	\choice 
		 {$(-2; 3;-1)$}
		 {\True $(2;-3; 1)$}
		 {$(2;-3;-1)$}
		{$(2; 3;  1)$}
	\loigiai{ 
		Trong không gian $Oxyz$, mặt cầu $(S)$ tâm $I(a; b; c)$  bán kính $R$ có phương trình là $(x-a)^2+(y-b)^2+(z-c)^2=R^2$.\\
		Do đó  mặt cầu $(S):(x-2)^2+(y+3)^2+(z-1)^2=16$ có tâm $I(2;-3; 1)$.
	 }
\end{ex}

% nguồn: Phan1_Lop_12_Hinh_hoc_Chuong_5
\begin{ex}%[2H5N3-2]%[ma:48ffaf7a]
	Trong không gian $Oxyz$, cho mặt cầu $(S)\colon (x-2)^2+(y+1)^2+(z-1)^2=9$. Tìm tọa  độ tâm $I$ và bán kính $R$ của mặt cầu $S$
	 \choice
		{$I(-2; 1;-1)$, $R=3$} 
		{$I(2;-1;  1)$, $R=9$}
		{\True $I(2;-1; 1)$,  $R=3$}
		{$I(-2; 1;-1)$,  $R=9$}
	\loigiai{Tâm $I$ và bán kính $R$ của mặt cầu $(S)\colon (x-2)^2+(y+1)^2+(z-1)^2=9$ là  $I(2;-1; 1)$, $R=3$.}
\end{ex}

% nguồn: Phan1_Lop_12_Hinh_hoc_Chuong_5
\begin{ex}%[2H5N3-2]%[ma:afef4a75]
	Trong không gian $Oxyz$, mặt cầu $(S)\colon (x - 2)^2 + (y + 1)^2 + (z - 3)^2 = 4$ có toạ độ tâm  $I$ là
	\choice 
		 {$(-2;1;-3)$}
		 {$(-2;1;3)$}
		{$(2;-1;-3)$} 
		{\True  $(2;-1;3)$}
	\loigiai{Tâm mặt cầu là $I(2;-1;3)$. 
	 }
\end{ex}

% nguồn: Phan1_Lop_12_Hinh_hoc_Chuong_5
\begin{ex}%[2H5N3-2]%[ma:b46a844d]
	Trong không gian $Oxyz$, cho mặt cầu $(S)\colon (x - 1)^2 + (y + 2)^2 +  (z - 1)^2 = 9$. Bán kính của $(S)$ là
	 \choice
		{$18$} 
		{$9$} 
		 {$6$}
		{\True  $3$}
	\loigiai{Phương trình mặt cầu dạng $(x - 1)^2 + (y  + 2)^2 + (z - 1)^2 = 9$. \\
	 So sánh với $(x - a)^2 + (y - b)^2 + (z - c)^2 = R^2$, suy ra bán kính $R = \sqrt{9} = 3$.
	 }
\end{ex}

% nguồn: Phan1_Lop_12_Hinh_hoc_Chuong_5
\begin{ex}%[2H5N3-2]%[ma:4b370456]
	Trong không gian $Oxyz$, phương trình của mặt cầu có tâm $I(1;0;-2)$ và bán kính $R=6$  là
	 \choice
		 {$(x-1)^{2}+y^{2}+(z-2)^{2}=6$}
		 {\True $(x-1)^{2}+y^{2}+(z+2)^{2}=36$}
		 {$(x-1)^{2}+y^{2}+(z+2)^{2}=6$}
		{$(x-1)^{2}+y^{2}+(z-2)^{2}=36$} 
	\loigiai{ 
		Phương trình mặt cầu  cần tìm là $(x-1)^2+y^2+(z+2)^2=36$.
	 }
\end{ex}

% nguồn: Phan1_Lop_12_Hinh_hoc_Chuong_5
\begin{ex}%[2H5N3-2]%[ma:cb80d6d4]
	Trong không gian $Oxyz$, cho mặt cầu $(S)\colon x^2+(y-2)^2+(z+1)^2=6$. Đường kính của $(S)$  bằng
	\choice 
		 {\True $2\sqrt{6}$}
		 {$\sqrt{6}$}
		{$6$} 
		 {$3$}
	 \loigiai{
		Mặt cầu $(S)$ có bán kính $R=\sqrt{6}$.\\ 
		Do  đó, đường kính của mặt cầu $(S)$ là $2R=2\sqrt{6}$.
	} 
\end{ex}

% nguồn: Phan1_Lop_12_Hinh_hoc_Chuong_5
\begin{ex}%[2H5N3-2]%[ma:845d5d5a]
	Trong không gian  $Oxyz$, phương trình
	$x^2+y^2+z^2-4x+2z-6y-2=0$ 
	là phương trình của mặt cầu  có tâm $I$. Xác định tọa độ tâm $I$.
	\choice 
		 {$I(2;-1;3)$}
		 {$I(2;3;-1)$}
		{$I(2;-3;1)$} 
		{\True $I(2;1;-3)$} 
	\loigiai{ 
		Phương trình mặt cầu 
		\begin{eqnarray*}
			&&x^2+y^2+z^2-4x+2z-6y-2=0\\ 
			&\Leftrightarrow& x^2+y^2+z^2-2\cdot2\cdot x-2\cdot 3\cdot  y-2\cdot (-1)\cdot z-2=0
		\end{eqnarray*}
		Vậy  tâm mặt cầu là $I(2;3;-1)$.
	} 
\end{ex}

% nguồn: Phan1_Lop_12_Hinh_hoc_Chuong_5
\begin{ex}%[2H5N3-2]%[ma:62ca30de]
	Trong không gian $Oxyz$, mặt cầu $(S)\colon x^2+y^2+z^2-6x+2y-15=0$ có bán kính $R$  bằng
	 \choice
		{\True $5$} 
		 {$\sqrt{5}$}
		 {$15$}
		 {$25$}
	\loigiai{ 
		Ta  có $a=\dfrac{-6}{-2}=3$, $b=\dfrac{2}{-2}=-1$, $c=\dfrac{0}{-2}=0$, $d=-15$.\\
		 Bán kính mặt cầu là
	 \[R=\sqrt{a^2+b^2+c^2-d}=\sqrt{3^2+(-1)^2+0^2+15}=5.\]}
\end{ex}

% nguồn: Phan1_Lop_12_Hinh_hoc_Chuong_5
\begin{ex}%[2H5N3-2]%[ma:d63c035e]
	Cho phương  trình mặt cầu $(S)\colon x^2+y^2+z^2=16$, bán kính $R$ của mặt cầu là
	 \choice
		{\True $R=4$} 
		{$R=16$} 
		 {$R=0$}
		{$R=2$} 
	\loigiai{ 
		Phương trình tổng quát của mặt cầu trong  không gian $Oxyz$ là $(x-a)^2+(y-b)^2+(z-c)^2=R^2$.\\
		Khi  đó $R^2=16 \Rightarrow R=4$.
	} 
\end{ex}

% nguồn: Phan1_Lop_12_Hinh_hoc_Chuong_5
\begin{ex}%[2H5N3-2]%[ma:f239eb06]
	Trong không gian $Oxyz$, cho mặt cầu $(S)$ có phương trình $(x+3)^{2}+(y-4)^{2}+z^{2}=81$. Điểm nào sau  đây là tâm của mặt cầu $(S)$?
	\choice 
		{$I_{1}(3;4;0)$} 
		{$I_{2}(3;-4;0)$} 
		 {\True $I_{4}(-3;4;0)$}
		{$I_{3}(-3;4;5)$} 
	 \loigiai{
		Tâm mặt cầu  $S$ là $I(-3;4;0)$.
	} 
\end{ex}

% nguồn: Phan1_Lop_12_Hinh_hoc_Chuong_5
\begin{ex}%[2H5N3-2]%[ma:62595e97]
	Trong không gian $Oxyz$, cho mặt cầu $(S)\colon x^2+y^2+z^2-2x+4y-6z-11=0$. Tọa độ tâm $I$  và bán kính $R$ của mặt cầu~là
	\choice 
		{$I(1;-2;3)$;  $R=25$}
		{\True  $I(1;-2;3)$; $R=5$}
		{$I(-1;2;-3)$; $R=5$} 
		{$I(-1;2;-3)$;  $R=25$}
	\loigiai{ 
	Mặt cầu $(S)\colon x^2+y^2+z^2-2x+4y-6z-11=0$ có tâm và bán kính~là  $I(1;-2;3)$; $R=5$.
	} 
\end{ex}

% nguồn: Phan1_Lop_12_Hinh_hoc_Chuong_5
\begin{ex}%[2H5N3-2]%[ma:75dc3913]
	 Trong không gian $Oxyz$, cho mặt cầu có phương trình
	\begin{align*}
		x^2 +y^2 +z^2  -2x +4y -6z-11=0.
	\end{align*}
	Tọa độ tâm  $I$ và bán kính $R$ của mặt cầu đã cho lần lượt~là
	\choice 
		{\True $I(1;-2;3)$,  $R=5$}
		 {$I(1;2;-3)$, $R=5$}
		{$I(-1;-2;3)$, $R=5$} 
		{$I(1;-2;-3)$, $R=5$} 
	 \loigiai{
		Phương trình $x^2 +y^2 +z^2 -2x +4y -6z-11=0$ có dạng  $x^2 +y^2 +z^2-2ax-2by-2cz+d=0$ với $a=1$, $b=-2$, $c=3$, $d=-11$.\\
		Ta có $a^2 +b^2  +c^2 -d=1+4+9+11=25>0$. Suy ra phương trình đã cho~là phương trình mặt cầu tâm $I(1;-2;3)$, bán kính $R=\sqrt{25}=5$.
	 }
\end{ex}

% nguồn: Phan1_Lop_12_Hinh_hoc_Chuong_5
\begin{ex}%[2H5N3-2]%[ma:9aee20c1]
	Trong không gian với hệ tọa  độ $Oxyz$, mặt cầu $(x+13)^2+(y-14)^2+(z-15)^2=4$ có bán kính bằng
	\choice 
		 {16}
		{\True 2} 
		 {256}
		{4} 
	\loigiai{Ta có $R^2=4 \Leftrightarrow  R=2$.}
\end{ex}

% nguồn: Phan1_Lop_12_Hinh_hoc_Chuong_5
\begin{ex}%[2H5N3-2]%[ma:b1dcb43f]
	Trong không gian với hệ tọa độ $Oxyz$, cho mặt cầu $(S)\colon (x-3)^2 + (y+1)^2 + (z-2)^2 = 8$. Đường kính của  mặt cầu này là
	\choice 
		 {\True $4\sqrt{2}$}
		{$4$} 
		{$2\sqrt{2}$} 
		 {$\sqrt{2}$}
	\loigiai{ 
		Mặt cầu $(S)$  có bán kính $R = \sqrt{8} = 2\sqrt{2}$. \\
		Đường kính  mặt cầu có $d = 2R = 2 \cdot 2\sqrt{2} = 4\sqrt{2}$.
	} 
\end{ex}

% nguồn: Phan1_Lop_12_Hinh_hoc_Chuong_5
\begin{ex}%[2H5N3-2]%[ma:1bd96231]
	Trong không gian $Oxyz$, tâm của mặt  cầu $(S)\colon (x+1)^2+(y-2)^2+(z-1)^2=9$ có tọa độ là
	\choice 
		{$(-1;  2;-1)$}
		{$(1;-2;  1)$}
		{$(1;-2;-1)$} 
		{\True $(-1; 2; 1)$} 
	\loigiai{ 
		Mặt cầu  $(S)$ có tâm $I(-1;2;1)$.
	 }
\end{ex}

% nguồn: Phan1_Lop_12_Hinh_hoc_Chuong_5
\begin{ex}%[2H5N3-2]%[ma:5fd690b6]
	Phương trình nào sau đây~là phương trình của một mặt cầu? 
	 \choice
		{$x^2+y^2+z^2-4  x-2 y+5=0$}
		{$x^2+y^2+z^2+2 x+6  y-2 z+15=0$}
		 {\True $x^2+y^2+z^2-4 x+1=0$}
		{$x^2+y^2+z^2-6 z+20=0$} 
	 \loigiai{
	Phương trình $x^2+y^2+z^2-2ax-2by-2cz+d=0$~là phương trình mặt cầu  khi $a^2+b^2+c^2-d>0$. \\
	Phương trình $x^2+y^2+z^2-4x+1=0$~là phương trình mặt cầu vì $2^2+0^2+0^2-1=3>0$. 
	} 
\end{ex}

% nguồn: Phan1_Lop_12_Hinh_hoc_Chuong_5
\begin{ex}%[2H5N3-2]%[ma:fd013e01]
	Trong không  gian $Oxyz$, cho mặt cầu có phương trình $(x-1)^2+(y+2)^2+(z-3)^2=4$. Tìm tọa độ tâm $I$ và bán kính $R$ của mặt cầu đó.
	 \choice
		{$I(-1;  2;-3)$; $R=2$}
		{$I(-1; 2;-3)$; $R=4$} 
		{\True $I(1;-2; 3)$;  $R=2$}
		{$I(1;-2;  3)$; $R=4$}
	 \loigiai{
		Mặt cầu  đã cho có tâm $I(1;-2; 3)$ và bán kính $R=2$.
	 }
\end{ex}

% nguồn: Phan1_Lop_12_Hinh_hoc_Chuong_5
\begin{ex}%[2H5N3-2]%[ma:47746fac]
	Trong không gian $Oxyz$, cho mặt cầu $(S)\colon x^2+y^2+z^2+2x-4y-2z-3=0$. Tọa độ tâm  $I$ của mặt cầu $(S)$ là
	\choice 
		 {\True$(-1;2;1)$}
		{$(1;-2;-1)$} 
		 {$(-2;4;2)$}
		 {$(2;-4;-2)$}
	\loigiai{ 
		Tọa độ  tâm $I(-1;2;1)$.
	} 
\end{ex}

% nguồn: Phan1_Lop_12_Hinh_hoc_Chuong_5
\begin{ex}%[2H5N3-2]%[ma:d60d62e2]
	Trong không gian với hệ trục tọa độ $Oxyz$, mặt cầu $(S)\colon(x+2)^2+(y-1)^2+z^2=4$ có tâm $I$ và bán kính $R$  lần lượt là
	 \choice
		{$I(2;-1; 0)$, $R=2$} 
		{$I(2;-1;  0)$, $R=4$}
		{$I(-2; 1; 0)$, $R=4$} 
		{\True  $I(-2; 1; 0)$, $R=2$}
	\loigiai{Mặt cầu có  tâm $I(-2;1; 0)$ và bán kính $R=2$.}
\end{ex}

% nguồn: Phan1_Lop_12_Hinh_hoc_Chuong_5
\begin{ex}%[2H5N3-2]%[ma:51d464c0]
	Trong không gian $Oxyz$, cho  mặt cầu $(S)$ có phương trình $(x-1)^2+(y+2)^2+(z-3)^2=4$. Tìm tọa độ tâm $I$ và bán kính $R$ của mặt cầu $(S)$.
	\choice 
		{$I=(-1;2;-3)$,  $R=2$}
		{$I=(-1;2;-3)$,  $R=4$}
		{\True  $I=(1;-2;3)$, $R=2$}
		 {$I=(1;-2;3)$, $R=4$}
	 \loigiai{
		 Tâm $I=(1;-2;3)$ và bán kính $R=\sqrt{4}=2$.
	 }
\end{ex}

% nguồn: Phan1_Lop_12_Hinh_hoc_Chuong_5
\begin{ex}%[2H5N3-2]%[ma:406facb6]
	Trong không  gian $Oxyz$, mặt cầu $(S) \colon x^2+y^2+z^2+4x-2y+8z-1=0$ có tọa độ tâm là
	\choice 
		{$(4;-2;8)$} 
		{$(2;-1;4)$} 
		{\True $(-2;1;-4)$} 
		 {$(2;-1;-4)$}
	 \loigiai{
		 Mặt cầu $(S)$ có tọa độ tâm là $I=(-2;1;-4)$.
	} 
\end{ex}

% nguồn: Phan1_Lop_12_Hinh_hoc_Chuong_5
\begin{ex}%[2H5N3-2]%[ma:51fbb2a7]
	Trong không gian $Oxyz$, cho mặt cầu $(S)\colon (x-1)^2+(y-2)^2+(z+3)^2=9$. Tọa độ tâm $I$ và bán  kính $R$ của ($S$) lần lượt là
	 \choice
		{$I(-1;2;-3)$; $R=3$} 
		 {$I(1;2;-3)$; $R=9$}
		{$I(1;-2;3)$; $R=9$} 
		{\True  $I(1;2;-3)$; $R=3$}
	\loigiai{ 
		Mặt cầu đã cho có tâm $I(1;2;-3)$ và  bán kính $R=3$.
	} 
\end{ex}

% nguồn: Phan1_Lop_12_Hinh_hoc_Chuong_5
\begin{ex}%[2H5N3-2]%[ma:6bf3abaf]
	Trong không gian $Oxyz$, xác định tâm $I$ và bán kính $R$ của mặt cầu $(S)$ có phương trình $(x-1)^2+(y-4)^2+(z+2)^2=9$ là 
	 \choice
		{\True  $I(1;4;-2),R = 3$}
		{$I(-1;-4;2),R  = 3$}
		{$I(1;4;-2),R = 9$} 
		{$I(-1;-4;2),R = 9$} 
	\loigiai{ 
		Mặt cầu $(S)$ có phương trình $(x-1)^2+(y-4)^2+(z+2)^2=9$ có tâm $I$ và bán  kính $R$ là $I (1;4;-2), R=3$.
	 }
\end{ex}

% nguồn: Phan1_Lop_12_Hinh_hoc_Chuong_5
\begin{ex}%[2H5N3-2]%[ma:85b89626]
	Trong không gian $Oxyz$, điểm nào  sau đây là tâm mặt cầu $(S)\colon x^{2}+y^{2}+z^{2}-2 x+4 y-1=0$.
	 \choice
		{$Q(1;  2;-1)$}
		{\True  $M(1;-2; 0)$}
		{$N(-1;-2; 0)$} 
		{$P(1;-2;-1)$} 
	\loigiai{ 
		Ta có $\heva{& -2a=-2\\& -2b=4\\& -2c=0}\Leftrightarrow\heva{& a=1\\& b=-2\\& c=0.}$\\ 
		 Vậy mặt cầu $(S)$ có tâm $M(1;-2; 0)$.
	} 
\end{ex}

% nguồn: Phan1_Lop_12_Hinh_hoc_Chuong_5
\begin{ex}%[2H5N3-2]%[ma:136c3c61]
	 Trong không gian $Oxyz$, cho mặt cầu $(S)\colon (x+1)^2+y^2+(z-3)^2=4$. Tọa độ tâm $I$ và bán kính $R$ của mặt cầu $(S)$ lần lượt là
	\choice 
		 {\True $I(-1;0;3)$, $R=2$}
		 {$I(-1;0;3)$, $R=4$}
		{$I(1;0;3)$,  $R=2$}
		 {$I(1;0;3)$, $R=4$}
	\loigiai{Tọa độ tâm $I$ và bán kính $R$ của mặt cầu $(S)$ lần lượt  là $I(-1;0;3)$, $R=2$.}
\end{ex}

% nguồn: Phan1_Lop_12_Hinh_hoc_Chuong_5
\begin{ex}%[2H5N3-2]%[ma:a14cac52]
	Trong không gian $Oxyz$, cho mặt cầu $(S) \colon x^2 +y^2 +z^2 -2x +2z -1=0$.  Tâm của $(S)$ có tọa~độ~là
	 \choice
		{$(-1;0;1)$} 
		{\True $(1;0;-1)$} 
		 {$(-1;1;0)$}
		{$(1;-1;0)$} 
	 \loigiai{
		Phương trình $x^2  +y^2 +z^2 -2x +2z -1=0$ có dạng
		$x^2  +y^2 +z^2-2ax-2by-2cz+d=0$ với $a=1$, $b=0$, $c=-1$, $d=-1$ nên tâm của mặt cầu~là $I(1;0;-1)$.
	} 
\end{ex}

% nguồn: Phan1_Lop_12_Hinh_hoc_Chuong_5
\begin{ex}%[2H5N3-2]%[ma:901bf501]
	Trong không gian $Oxyz$, phương trình $x^2+y^2+z^2-4x+2z-6y-2=0$ là phương trình của mặt cầu có tâm $I$. Xác định  tọa độ tâm $I$
	\choice 
		 {$I(2;-1;3)$}
		 {\True $I(2;3;-1)$}
		 {$I(2;-3;1)$}
		{$I(2;1;-3)$} 
	\loigiai{ 
		Mặt cầu tâm $I(a;b;c)$, bán kính  $R$ có phương trình $x^2+y^2+z^2-2ax-2by-2cz+d=0$.\\
		 $\Rightarrow\heva{&-2a=-4\\&-2b=-6\\&-2c=2} \Leftrightarrow\heva{&a=2\\&b=3\\&c=-1.}$\\
		Vậy tâm của mặt cầu  là $I(2;3;-1)$.
	} 
\end{ex}

% nguồn: Phan1_Lop_12_Hinh_hoc_Chuong_5
\begin{ex}%[2H5N3-2]%[ma:bcbe155e]
	Trong không gian $Oxyz$, cho mặt cầu $(S)$ có phương trình  $x^2+y^2+z^2-2x-4y-6z=0$. Trong ba điểm có tọa độ lần lượt là $\left(0;0;0\right)$, $\left(1;2;3\right)$, $\left(2;0;6\right)$ thì có bao nhiêu điểm nằm trên mặt cầu $(S)$ ?
	 \choice
		{$0$} 
		{$3$} 
		{$1$} 
		 {\True $2$}
	 \loigiai{
		Thay tọa độ ba điểm vào phương trình mặt  cầu, ta có\\
		 $0^2+0^2+0^2-2\cdot0-4\cdot0-6\cdot0=0$ (Đúng) $\Rightarrow O\left(0;0;0\right)\in(S)$.\\
		$1^2+2^2+3^2-2\cdot1-4\cdot2-6\cdot3=0$ (Sai)  $\Rightarrow A\left(1;2;3\right)\notin(S)$.\\
		$2^2+0^2+6^2-2\cdot2-4\cdot0-6\cdot6=0$ (Đúng) $\Rightarrow B\left(2;0;6\right)\in(S)$.\\ 
	Vậy trong ba điểm đề cho, có hai điểm nằm trên  mặt cầu $(S)$.}
\end{ex}

% nguồn: Phan1_Lop_12_Hinh_hoc_Chuong_5
\begin{ex}%[2H5N3-2]%[ma:78af030a]
	Trong không  gian $Oxyz$, cho mặt cầu $(S)$ có phương trình $(x-3)^2+(y-1)^2+(z+2)^2=9$. Tọa độ tâm $I$ và bán kính $R$ của mặt cầu là
	 \choice
		{$I(3;1;-2)$, $R=9$} 
		{\True  $I(3;1;-2)$, $R=3$}
		 {$I(-3;-1;2)$, $R=3$}
		{$I(-3;-1;2)$,  $R=9$}
	\loigiai{Mặt cầu $(S)$ có phương trình  $(x-3)^2+(y-1)^2+(z+2)^2=9$ có tâm $I(3;1;-2)$ và $R=3$.}
\end{ex}

% nguồn: Phan1_Lop_12_Hinh_hoc_Chuong_5
\begin{ex}%[2H5H3-2]%[ma:ecccefde]
	Trong không gian $Oxyz$, phương trình nào trong các phương trình sau là phương trình của  một mặt cầu?
	\choice 
		{$(S)\colon  x^2+y^2-z^2-4x+2y+2z-3=0$}
		 {\True $(S)\colon x^2+y^2+z^2-4x+2y+2z+2=0$}
		{$(S)\colon  x^2+2y^2+z^2-4x+2y+2z-10=0$}
		 {$(S)\colon x^2+y^2-4x+2y+2z+8=0$}
	\loigiai{ 
		Phương trình $(S)\colon x^2+y^2+z^2-4x+2y+2z+2=0$ là phương  trình mặt cầu tâm $I(2;-1;-1)$ và bán kính $R=\sqrt{2^2+(-1)^2+(-1)^2-2}=2$.
	} 
\end{ex}

% nguồn: Phan1_Lop_12_Hinh_hoc_Chuong_5
\begin{ex}%[2H5H3-2]%[ma:d316dff9]
	Cho  mặt cầu $(S)\colon (x+1)^2+(y-2)^2+(z-1)^2=9$.
	Toạ tâm $I$ và bán kính $R$ của  $(S)$ là
	\choice 
		{\True $I(-1;  2; 1)$ và $R=3$}
		{$I(1;-2;-1)$ và  $R=3$}
		{$I(-1;  2; 1)$ và $R=9$}
		{$I(1;-2;-1)$  và $R=9$}
	 \loigiai{
		Mặt  cầu $(S)\colon (x+1)^2+(y-2)^2+(z-1)^2=9$
		có  tâm $I(-1; 2; 1)$ và bán kính $R=3$.
	} 
\end{ex}

% nguồn: Phan1_Lop_12_Hinh_hoc_Chuong_5
\begin{ex}%[2H5H3-2]%[ma:40edcaf9]
	Trong không gian $Oxyz$,  mặt cầu $(S)\colon(x-2)^2+(y-3)^2+(z-4)^2=9$ có bán kính bằng
	\choice 
		{$81$} 
		 {$9$}
		{\True $3$} 
		 {$6$}
	 \loigiai{
		Mặt cầu $(S)\colon(x-2)^2+(y-3)^2+(z-4)^2=9$  có bán kính $R=\sqrt{9}=3$.
	} 
\end{ex}

% nguồn: Phan1_Lop_12_Hinh_hoc_Chuong_5
\begin{ex}%[2H5H3-2]%[ma:98298faa]
	Trong không gian $Oxyz$, cho mặt cầu $(S)\colon x^2+y^2+z^2-2x+4y-4z-25=0$. Tìm tọa độ tâm $I$ và bán kính $R$ của mặt cầu $(S)$? 
	 \choice
		{$I(1;-2;2)$, $R=6$}{$I(-1;2;-2)$, $R=5$}{\True $I(1;-2;2)$, $R=\sqrt{34}$}{$I(-2;4;-4$; $R=\sqrt{29}$} 
	\loigiai{ 
		Mặt cầu $(S)\colon x^2+y^2+z^2-2x+4y-4z-25=0\Rightarrow  a=1$, $b=-2$; $c=2$, $d=-25$.\\
		Do đó  mặt cầu $(S)$ có tâm là $I(1;-2;2)$, bán kính $R=\sqrt{1^2+(-2)^2+2^2+25}=\sqrt{34}$.
	} 
\end{ex}

% nguồn: Phan1_Lop_12_Hinh_hoc_Chuong_5
\begin{ex}%[2H5H3-2]%[ma:546faeef]
	Trong không gian với hệ trục tọa độ $Oxyz$, cho mặt cầu $(S)$ có phương trình $x^2+y^2+z^2-2x-4y+6z-2=0$. Đường kính của mặt cầu $(S)$ bằng 
	\choice 
		{\True $8$} 
		{$16$} 
		{$2\sqrt{14}$} 
		{$4$} 
	 \loigiai{
		Ta có $R=\sqrt{1+4+9+2}=4$ suy ra đường  kính của mặt cầu $(S)$ bằng $8$.
	 }
\end{ex}

% nguồn: Phan1_Lop_12_Hinh_hoc_Chuong_5
\begin{ex}%[2H5H3-2]%[ma:2b15c7fb]
	Xác định toạ độ tâm và bán kính của  mặt cầu $(S)\colon x^2+y^2+z^2+8x-6y+2z-10=0$.
	\choice 
		{\True Tâm $I(-4;3;-1)$ và bán kính $R=6$} 
		{Tâm $I(-4;3;-1)$ và bán kính  $R=36$}
		{Tâm $I(4;-3;1)$ và bán kính  $R=6$}
		{Tâm $I(4;-3;1)$ và bán kính $R=36$} 
	\loigiai{Mặt  cầu $(S)\colon x^2+y^2+z^2+8x-6y+2z-10=0$ có
		 $\heva{&a=\dfrac{8}{-2}=-4\\&b=\dfrac{-6}{-2}=3\\&c=\dfrac{2}{-2}=-1\\&d=-10.}$\\
		 Vậy mặt cầu đã cho có tâm $I(-4;3;-1)$ và bán kính $R=\sqrt{(-4)^2+3^2+(-1)^2+10}=6$.
	} 
\end{ex}

% nguồn: Phan1_Lop_12_Hinh_hoc_Chuong_5
\begin{ex}%[2H5H3-2]%[ma:44d455b8]
	Trong không gian $Oxyz$, mặt cầu $(S)$ có tâm $M(1;2;4)$ tiếp xúc mặt phẳng $(P)\colon2x+y-2z+1=0$ có  bán kính bằng
	 \choice
		{$\dfrac{4}{3}$} 
		{\True $1$} 
		{$\dfrac{3}{4}$} 
		{$2$} 
	\loigiai{Do mặt cầu $(S)$  tiếp xúc với mặt phẳng $(P)$ nên bán kính của $(S)$ là
	 \[R=\mathrm{d}(I,(P))=\dfrac{|2+2-8+1|}{\sqrt{4+1+4}}=1.\]}
\end{ex}

% nguồn: Phan1_Lop_12_Hinh_hoc_Chuong_5
\begin{ex}%[2H5H3-2]%[ma:40286ce5]
	Trong  không gian $Oxyz$, gọi $I$ là tâm của mặt cầu $(S)\colon x^2+y^2+z^2+2x-4z-1=0$. Độ dài đoạn $OI$ (với $O$ là gốc tọa độ) bằng
	 \choice
		{$5$} 
		 {\True $\sqrt{5}$}
		 {$\sqrt{6}$}
		{$6$} 
	 \loigiai{
	 Ta có $I\left(-1;0;2\right)$. Suy ra $OI=\sqrt{\left(-1\right)^2+0^2+2^2}=\sqrt{5}$.}
\end{ex}

% nguồn: Phan1_Lop_12_Hinh_hoc_Chuong_5
\begin{ex}%[2H5H3-2]%[ma:0f947c7e]
	Trong  không gian $Oxyz$, cho mặt cầu $(S)\colon x^2+y^2+z^2-2x+4y+2z-10=0$. Tâm và bán kính của mặt cầu $(S)$ là
	 \choice
		{$I(-2;4;2)$,  $R=4$}
		{\True  $I(1;-2;-1)$, $R=4$}
		 {$I(2;-4;-2)$, $R=\sqrt{10}$}
		{$I(-1;2;1)$, $R=\sqrt{34}$} 
	\loigiai{Mặt cầu $(S)\colon x^2+y^2+z^2-2x+4y+2z-10=0$ có tâm $I(1;-2;-1)$ và bán kính\\  $R=\sqrt{1^2+(-2)^2+(-1)^2-(-10)}=\sqrt{16}=4$.}
\end{ex}

% nguồn: Phan1_Lop_12_Hinh_hoc_Chuong_5
\begin{ex}%[2H5H3-2]%[ma:007cde09]
	Trong không gian $Oxyz$, cho mặt cầu $\left(S\right)\colon\left(x+1\right)^2+\left(y-3\right)^2+\left(z+2\right)^2=4$. Tìm tọa  độ tâm $I$ và bán kính $R$ của mặt cầu $\left(S\right)$.
	\choice 
		{$I\left(1;-3;2\right)$, $R=4$} 
		{$I\left(1;-3;2\right)$, $R=2$} 
		 {\True $I\left(-1;3;-2\right)$, $R=2$}
		 {$I\left(-1;3;-2\right)$, $R=4$}
	 \loigiai{
		Mặt cầu $(S)\colon(x+1)^2+(y-3)^2+(z+2)^2=4$ có tâm $I(-1;3;-2)$ và bán  kính $R=2$.
	 }
\end{ex}

% nguồn: Phan1_Lop_12_Hinh_hoc_Chuong_5
\begin{ex}%[2H5V3-2]%[ma:bd58a535]
	Trong không gian $Oxyz$, cho điểm $M$ thuộc mặt cầu $(S)\colon\left(x-3\right)^2+\left(y-3\right)^2+\left(z-2\right)^2=9$ và ba điểm $A\left(1;0;0\right)$, $B\left(2;1;3\right)$, $C\left(0;2;-3\right)$. Biết rằng tập hợp các điểm $M$  thỏa mãn biểu thức $MA^2+2\cdot\overrightarrow{MB}\cdot\overrightarrow{MC}=8$ là đường tròn cố định, bán kính $r$ của đường tròn này là
	 \choice
		{\True $r=\sqrt{6}$} 
		{$r=3$} 
		 {$r=6$}
		 {$r=\sqrt{3}$}
	\loigiai{ 
		Gọi $M\left(x;y;z\right)$ khi đó $M{A^2}+2\cdot\overrightarrow{MB}\cdot\overrightarrow{MC}=8$\\ 
		 $\Leftrightarrow\left(x-1\right)^2+y^2+z^2+2\cdot\left[\left(2-x\right)\cdot\left(-x\right)+\left(1-y\right)\cdot\left(2-y\right)+\left(3-z\right)\cdot\left(-3-z\right)\right]=8$\\
		$\Leftrightarrow{x^2}+y^2+z^2-2x-2y-7=0\Leftrightarrow{\left(x-1\right)^2}+\left(y-1\right)^2+z^2=9$ $\Rightarrow M\in\left(S_1\right)$.\\ 
		Mặt cầu $(S)$ có tâm $I(3;3;2)$, bán kính $R=3$. Mặt cầu $\left(S_1\right)$ có tâm $I_1(1;1;0)$,  bán kính $R_1=3$.\\
		 Ta có $II_1=2\sqrt{3}<R+R_1\Rightarrow (S)$ cắt mặt cầu $\left(S_1\right)$ theo giao tuyến là một đường tròn.\\
		 Đường tròn này nằm trên mặt phẳng\\
		 $(P)\colon\left(x-3\right)^2+\left(y-3\right)^2+\left(z-2\right)^2-\left[\left(x-1\right)^2+\left(y-1\right)^2+z^2\right]=0\Leftrightarrow x+y+z-5=0$.\\
	Khi  đó $\mathrm{d}\left(I;(P)\right)=\sqrt{3}\Rightarrow r=\sqrt{R^2-\mathrm{d}^2\left(I;(P)\right)}=\sqrt{6}$.}
\end{ex}

% nguồn: Phan1_Lop_12_Hinh_hoc_Chuong_5
\begin{ex}%[2H5V3-2]%[ma:1ab7d913]
	Trong không gian $Oxyz$, cho hai  điểm $A\left(1;5;2\right)$ và $B\left(5;13;10\right)$. Có bao nhiêu điểm $I\left(a;b;c\right)$ với $a$, $b$, $c$ là các số nguyên sao cho có mặt cầu tâm $I$ đi qua $A$, $B$ và tiếp xúc với mặt phẳng $\left(Oxy\right)$?
	\choice 
		 {$10$}
		 {$6$}
		{$8$} 
		{\True  $4$}
	 \loigiai{
		 Ta có $\overrightarrow{AB}=\left(4;8;8\right)\Rightarrow AB\colon\heva{
		& x=1+t\\ 
		& y=5+2t\\ 
		 & z=2+2t.}$\\
		Gọi $M$ là trung điểm của $AB\Rightarrow  M\left(3;9;6\right)\Rightarrow$ Mặt phẳng trung trực của $AB$ là\\
		$\left(\alpha\right)\colon x-3+2\left(y-9\right)+2\left(z-6\right)=0\Leftrightarrow x+2y+2z-33=0$.\\ 
		$I\in\left(\alpha\right)\Rightarrow a+2b+2c-33=0\Leftrightarrow a=33-2b-2c\Rightarrow a$ là số  nguyên lẻ.\\
		Gọi  $J=AB\cap\left(Oxy\right)\Rightarrow J\left(0;3;0\right)\Rightarrow JA=3$, $JB=15$.\\
		Gọi  $C$ là tiếp điểm của mặt cầu và $\left(Oxy\right)\Rightarrow C\left(a;b;0\right)$.\\
		Ta có $JA\cdot JB=JC^2\Rightarrow JC^2=45\Rightarrow C$ thuộc  đường tròn tâm $\left(J,3\sqrt{5}\right)$.\\
		Xét trong  mặt phẳng $\left(Oxy\right)$, phương trình của $\left(J,3\sqrt{5}\right)\colon a^2+\left(b-3\right)^2=45\Rightarrow a^2\le 45$.\\
		Do $a$, $b$, $c\in\mathbb{Z}$ và $a$  lẻ nên ta có $a^2=1\Rightarrow{\left(b-3\right)^2}=44$ (loại);\\
		$a^2=9\Rightarrow{\left(b-3\right)^2}=36$ (thỏa  mãn); $a^2=25\Rightarrow{\left(b-3\right)^2}=20$ (loại).\\
	Vậy có  $4$ bộ $\left(a;b;c\right)$ thỏa mãn.}
\end{ex}

% nguồn: Phan1_Lop_12_Hinh_hoc_Chuong_5
\begin{ex}%[2H5V3-2]%[ma:ac2fe9b5]
	Trong không gian $Oxyz$, cho điểm $M(6;5;4)$ và mặt  cầu $(S)$ đi qua $M$, tiếp xúc với mặt phẳng $(Oxy)$. Gọi $R$ là bán kính của mặt cầu $(S)$, giá trị nhỏ nhất của $R$ là
	\choice 
		{$3$} 
		 {$\sqrt{61}$}
		{$4$} 
		{\True  $2$}
	\loigiai{ 
		\begin{center}
			\begin{tikzpicture}[line join = round,line cap = round,>=stealth,font = \small,scale = 1]
					\def  \a{2.5}
					 \def \b{0.785*\a}
					\path (0:0) coordinate (O) 
						(0:8.3) coordinate (Y) 
						(40:2)  coordinate (X)
						(35:5) coordinate  (I)
						 ++(-90:\a) coordinate (A);
					\path (I) +(35:\a)  coordinate (M);
					\path (A) +(0:\a) coordinate (A'); 
					\path  (M) +(-90:2*\a) coordinate (M');
					\path (intersection of A--A' and M--M') coordinate (B); 
					 \path ($(M)!.5!(B)$) coordinate (I_0);
					\draw [name path=c] (I) circle  (\a);
					\draw[gray] let \p1=($(I_0)-(B)$)  in (I_0) circle ({veclen (\x1,\y1)});
					\path  ($(A)!2!(I)$) coordinate (N);
					 \draw (X)--(O)--(Y);
					 \path[name path=d] (I_0)--(B);
					\draw[thin] pic[draw,angle radius = 12pt,"$(Oxy)$",angle eccentricity = 2.5]{angle  = Y--O--X};
					\path [name intersections = {of =  d and c}]
						(intersection-1) coordinate (E); 
					\draw[dashed] (N)--(A) 
						(M)--(E); 
					\draw (E)--(B); 
					\foreach \x  in {I,A,B,N,E} \fill (\x) circle (1.5pt);
					\foreach \x/\g in  {I/180,I_0/220,M/45}
					\fill (\x)  circle (1.5pt)
						+(\g:3mm)  node {$\x$};
					\path (I) +(0:\a)  coordinate (C);
					\draw[dashed] (C) arc  (0:180:{\a} and {0.2*\a});
					\draw (C) arc (0:-180:{\a}  and {0.2*\a});
					\path (I_0) +(0:\b)  coordinate (D);
					\draw[dashed,gray]  (D) arc (0:180:{\b} and {0.2*\b});
					\draw[gray] (D) arc  (0:-180:{\b} and {0.2*\b});
					\foreach \x/\y/\z in  {I/A/B,I_0/B/A'}{
						\draw[thin] pic[draw,angle radius = 6pt]{right angle = \x--\y--\z}; 
					 }
				\end{tikzpicture}
		\end{center}
		Gọi $\mathrm{d}$ là khoảng cách từ $M$ đến mặt phẳng  $(Oxy) \Rightarrow \mathrm{d}=4$.\\
		Vì mặt cầu $(S)$ tiếp xúc $(Oxy)$ và luôn đi qua $M$ nên ta có đánh  giá $2R \geq \mathrm{d} \Leftrightarrow 2R \geq 4 \Leftrightarrow R \geq 2$.\\
		Vậy  giá trị nhỏ nhất của $R$ là $2$.
	 }
\end{ex}

% nguồn: Phan1_Lop_12_Hinh_hoc_Chuong_5
\begin{ex}%[2H5V3-2]%[ma:23389a5a]
	Trong không gian với hệ tọa độ $Oxyz$, cho hai mặt cầu $\left(S_1\right)$, $\left(S_2\right)$ thỏa $\left(S_1\right)$ có tâm $I\left(0;0;\dfrac{21}{2}\right)$, bán kính $r_1=6$ và $\left(S_2\right)$ có tâm $J\left(0;0;1\right)$, bán kính $r_2=\dfrac{9}{2}$. Hỏi có bao nhiêu điểm $M\left(x;y;z\right)$ với $x$, $y$, $z$ nguyên thuộc phần  giao của hai khối cầu?
	\choice 
		{$11$} 
		{\True  $13$}
		{$9$} 
		{$7$} 
	 \loigiai{
		Ta có phương  trình mặt cầu $\left(S_1\right)\colon x^2+y^2+\left(z-\dfrac{21}{2}\right)^2=36$.\\
		 Và phương trình mặt cầu $\left(S_2\right)\colon x^2+y^2+\left(z-1\right)^2=\dfrac{81}{4}$.\\
		Điểm $M\left(x;y;z\right)$ thuộc giao của hai khối cầu  $\left(S_1\right)$, $\left(S_2\right)$ nên toạ độ điểm $M\left(x;y;z\right)$ là nghiệm của hệ bất phương trình\\
		$\heva{&{x^2}+y^2+\left(z-\dfrac{21}{2}\right)^2\le 36\\ 
		 &{x^2}+y^2+\left(z-1\right)^2\le\dfrac{81}{4}}
		 \Leftrightarrow\heva{&x^2+y^2+\left(z-1\right)^2\le\dfrac{81}{4}\\
		&  z=5.}$\\
		Từ  đó suy ra $x^2+y^2+\left(5-1\right)^2\le\dfrac{81}{4}\Leftrightarrow{x^2}+y^2\le\dfrac{17}{4}$.\\
		Do $x$, $y\in\mathbb{Z}$ và $x^2+y^2\le\dfrac{17}{4}$ suy ra  $\heva{& x=0\\
		& y=\pm 1}$;  $\heva{
		& x=\pm  1\\
		&  y=0}$;$\heva{
		& x=\pm  1\\
		&  y=\pm 1}$;$\heva{
		& x=0\\ 
		&  y=\pm 2}$;$\heva{
		& x=\pm  2\\
		&  y=0}$; $\heva{
		 & x=0\\
		 & y=0.}$\\
		Vậy có $13$ điểm $M\left(x;y;z\right)$ với $x$, $y$, $z$ nguyên thuộc phần giao của hai khối cầu. 
	} 
\end{ex}

% nguồn: Phan1_Lop_12_Hinh_hoc_Chuong_5
\begin{ex}%[2H5V3-2]%[ma:40a88989]
	Trong không gian $Oxyz$, cho mặt cầu $(S)\colon\left(x-1\right)^2+\left(y+2\right)^2+\left(z-4\right)^2=4$. Có bao nhiêu điểm $P$ thuộc $(S)$ mà tiếp diện của $(S)$ tại $P$  cắt các trục $Ox$, $Oz$ tương ứng tại các điểm $E\left(a;0;0\right)$, $F\left(0;0;b\right)$ sao cho $a$, $b$ là các số nguyên dương và $\widehat{EPF}=90^{\circ}$?
	\choice 
		{$1$} 
		{$2$} 
		 {\True $4$}
		{$3$} 
	\loigiai{ 
		Mặt cầu $(S)$ có tâm $I\left(1;-2;4\right)$ và bán kính $R=2$.\\ 
		Ta có $\mathrm{d}\left(I,\left(Oxz\right)\right)=\dfrac{\left|-2\right|}{1}=2=R$ nên  $(S)$ tiếp xúc với $\left(Oxz\right)$ tại $H\left(1;0;4\right)$.\\
		Dễ thấy $EH=EP$, $FH=FP$ (do tính chất các tiếp tuyến kẻ từ $1$  điểm đến mặt cầu).\\
		Suy ra $\triangle EPF=\triangle EHF$ (c  – c – c) nên $\widehat{EHF}=\widehat{EPF}=90^{\circ}$ ($2$ góc tương ứng) hay $\overrightarrow{HE}\cdot\overrightarrow{HF}=0$.\\
		Khi đó  $\left(a-1;0;-4\right)\cdot\left(-1;0;b-4\right)=0\Leftrightarrow-a+1-4b+16=0\Leftrightarrow a+4b-17=0$.\\
		Vì $a$ nguyên dương  nên $a=17-4b\ge 0\Rightarrow b\le\dfrac{17}{4}\Rightarrow b\in\left\{1;2;3;4\right\}$.\\
		Các cặp $\left(a;b\right)$ thỏa mãn  là $\left(13;1\right)$, $\left(9;2\right)$, $\left(5;3\right)$, $\left(1;4\right)$.
	} 
\end{ex}


%%%%%%%%%%%%%%%% ĐÚNG/SAI %%%%%%%%%%%%%%%%

% nguồn: Phan1_Lop_12_Hinh_hoc_Chuong_5
\begin{ex}%[2H5H3-2]%[ma:41a8f200]
	Trong không gian với hệ trục tọa độ $Oxyz$, cho đường thẳng $d\colon x=\dfrac{y-1}{2}=\dfrac{z-2}{-1}$ và mặt phẳng $(P)\colon 3x+y-z-5=0$. Mặt cầu $(S)$ có tâm $I$ thuộc đường thẳng $d$  và cắt mặt phẳng $(P)$ theo giao tuyến là đường tròn lớn nhất có bán kính $r=5$.
	\choiceTF 
		{\True Mặt phẳng $(P)\colon  3x+y-z-5=0$ có vectơ pháp tuyến $\overrightarrow{n}=\left(3;1;-1\right)$}
		{Tọa độ tổng quát của tâm  $I$ là $\left(t;-1+2t;-2-t\right)$}
		 {$\mathrm{d}\left(I,(P)\right)=3$}
		{\True Mặt cầu $(S)$ có  phương trình là $\left(x-1\right)^2+\left(y-3\right)^2+\left(z-1\right)^2=25$}
	\loigiai{\begin{itemchoice}
		\itemch $(P)\colon3x+y-z-5=0$  có vectơ pháp tuyến là $\overrightarrow{n}=\left(3;1;-1\right)$.
		\itemch Ta  có $d\colon x=\dfrac{y-1}{2}=\dfrac{z-2}{-1}\Rightarrow d\colon\heva{& x=t\\
		 &y=1+2t\\
		&z=2-t,}$ $t\in\mathbb{R}$.\\ 
		Vì $I\in d\Rightarrow  I\left(t;1+2t;2-t\right)$.
		\itemch Do giao tuyến  của $(S)$ và $(P)$ là đường tròn lớn nên $I\in(P)\Rightarrow \mathrm{d}\left(I,(P)\right)=0$.
		\itemch Ta  có $I\in d\Rightarrow I\left(t;1+2t;2-t\right)$.\\
		Theo giả thiết $I=d\cap(P)$ nên tọa độ điểm $I$  thỏa mãn phương trình mặt phẳng $(P)$.\\
		Thay tọa độ điểm $I$ vào  $(P)$ ta có\\
		$3t+\left(1+2t\right)-\left(2-t\right)-5=0\Leftrightarrow 6t-6=0\Leftrightarrow t=1\Rightarrow I\left(1;3;1\right)$.\\ 
		Vì $(S)\cap(P)$ theo giao  tuyến là đường tròn lớn có bán kính $r=5$ nên ta có bán kính mặt cầu $(S)$ là $R=r=5$.\\
		 Vậy mặt cầu có phương trình là $\left(x-1\right)^2+\left(y-3\right)^2+\left(z-1\right)^2=25$.
	\end{itemchoice}
 }
\end{ex}

% nguồn: Phan1_Lop_12_Hinh_hoc_Chuong_5
\begin{ex}%[2H5H3-2]%[ma:667cc963]
	 Trong không gian $Oxyz$, cho mặt cầu $(S)\colon (x+1)^2+(y-1)^2+(z+2)^2=9$ có tâm $I$ và bán kính $R$.
	 \choiceTF
		{Phương  trình mặt cầu $(S)$ có tâm $I(1;-1;2)$, bán kính $R=3$}
		{\True Điểm $A(0;2;-3)$ nằm trong mặt  cầu}
		{Điểm $J(1;2;3)$ nằm ngoài mặt cầu và khoảng cách từ tâm $I$ đến điểm $J$ bằng  $\sqrt{10}$}
		{Khoảng  cách từ tâm $I$ đến tâm mặt cầu $(S')\colon x^2+y^2+(z-1)^2=3$ bằng $\sqrt{2}$}
	 \loigiai{
		\begin{itemchoice}
			\itemch Tâm mặt cầu là $I(-1;1;-2)$ và  bán kính $R=3$.
			\itemch Vì $IA=\sqrt{(0+1)^2+(2-1)^2+(-3+2)^2}=\sqrt{1+1+1}=\sqrt{3}<3$ nên điểm $A$ nằm trong mặt  cầu.
			\itemch Vì $IJ=\sqrt{(1+1)^2+(2-1)^2+(3+2)^2}=\sqrt{4+1+25}=\sqrt{30}>3$, nên $J$ nằm  ngoài mặt cầu.
			\itemch Tâm  mặt cầu $(S')$ là $I'(0;0;1)$.\\ $$II'=\sqrt{(-1-0)^2+(1-0)^2+(-2-1)^2}=\sqrt{11}\ne\sqrt{2}.$$
		\end{itemchoice}
	} 
\end{ex}

% nguồn: Phan1_Lop_12_Hinh_hoc_Chuong_5
\begin{ex}%[2H5H3-2]%[ma:a2b0f1c6]
	Trong không gian $Oxyz$, cho mặt cầu $(S)\colon x^2+y^2+z^2-2x+4z+1=0$. Xét tính đúng sai của  các khẳng định sau
	\choiceTF 
		{\True Tọa độ tâm  $I$ của mặt cầu là $(1; 0;-2)$}
		{\True  Diện tích mặt cầu bằng $16\pi$ (đvdt)}
		{Điểm $A(1; 2; 3)$ nằm  trong mặt cầu}
		{\True Số điểm chung của đường thẳng $\Delta\colon \dfrac{x}{2}=\dfrac{y-1}{1}=\dfrac{z-2}{-1}$ và mặt cầu $(S)$  bằng $0$}
	\loigiai{ 
		\begin{itemchoice}
			\itemch Tọa độ tâm  $I$ của mặt cầu là $\heva{&x_I=-2:-2=1\\&y_I=0:-2=0\\&z_I=4:-2=-2} \Rightarrow I(1; 0;-2)$.
			\itemch $R=\sqrt{1^2+0^2+(-2)^2-1}=2\Rightarrow  S_{m c}=4\pi R^2=4\pi \cdot 2^2=16\pi$ (đvdt).
			\itemch  $\overrightarrow{IA}=(0; 2; 5) \Rightarrow\big|\overrightarrow{IA}\big|=\sqrt{0^2+2^2+5^2}=\sqrt{29} > 2\Rightarrow$ điểm $A$ nằm ngoài mặt cầu.
			\itemch Đường thẳng $\Delta$ đi qua $M(0; 1; 2)$ và nhận $\vec{u}=(2;  1;-2)$ làm một vectơ chỉ phương.\\
			 $\mathrm{d}(I, \Delta)=\dfrac{\bigg|\left[\vec{u}, \overrightarrow{MI}\right]\bigg|}{|\vec{u}|}=\dfrac{\sqrt{498}}{6} > R\Rightarrow \Delta$ không cắt $(S) \Rightarrow$ số điểm chung của đường thẳng $\Delta$ và mặt cầu $(S)$ bằng $0$.
		\end{itemchoice}
	} 
\end{ex}

% nguồn: Phan1_Lop_12_Hinh_hoc_Chuong_5
\begin{ex}%[2H5H3-2]%[ma:50abef89]
	Trong không gian $Oxyz$, cho mặt cầu $(S)\colon x^2+y^2+z^2-2x+4y+1=0$ có  tâm $I$ và bán kính $R$.
	\choiceTF 
		{\True Mặt  cầu $(S)$ có tâm $I\left(1;-2;0\right)$ và bán kính $R=2$}
		{Bán kính của mặt cầu $(S)$ là đoạn $IM$ với điểm  $M\left(1;1;2\right)$}
		{Mặt cầu  $(S)$ có đường kính $AB$ với $A\left(0;1;-2\right)$ và $B\left(2;-1;-4\right)$}
		{Mặt  cầu $(S)$ tiếp xúc với mặt phẳng $(P)\colon x+y-z-2=0$}
	 \loigiai{
		\begin{itemchoice}
			\itemch Mặt cầu $(S)\colon x^2+y^2+z^2-2x+4y+1=0$ có $a=1$, $b=-2$, $c=0$, $d=1$.\\ 
			$\Rightarrow $ tâm $I\left(1;-2;0\right)$  và bán kính $R=\sqrt{a^2+b^2+c^2-d}=2$.
			\itemch  Xét $M\left(1;1;2\right)\colon 1^2+1^2+2^2-2\cdot1+4\cdot1+1=9$ do đó $M$ không nằm trên $(S)$.\\
			 Vậy bán kính mặt cầu $(S)$ không là đoạn $IM$.
			\itemch Mặt cầu $(S)$ nhận $AB$ làm đường kính thì $\dfrac{AB}{2}=R$ và tâm $I$ là trung  điểm của $AB$.\\
			Ta  có $\overrightarrow{AB}=\left(2;-2;-2\right)\Rightarrow\left|\overrightarrow{AB}\right|=\sqrt{2^2+\left(-2\right)^2+\left(-2\right)^2}=2\sqrt{3}\Rightarrow\dfrac{AB}{2}=\sqrt{3}\ne R=2$.
			\itemch Mặt cầu  $(S)$ tiếp xúc với mặt phẳng $(P)\colon x+y-z-2=0\Leftrightarrow R=\mathrm{d}\left(I,(P)\right)$.\\
			Ta có $\mathrm{d}\left(I,(P)\right)=\dfrac{\left| 1+\left(-2\right)-0-2\right|}{\sqrt{1^2+1^2+\left(-1\right)^2}}=\sqrt{3}\ne  R=2$.
		\end{itemchoice}
	 }
\end{ex}

% nguồn: Phan1_Lop_12_Hinh_hoc_Chuong_5
\begin{ex}%[2H5H3-2]%[ma:eef3c5c6]
	Trong không gian  $Oxyz$, cho mặt cầu $(C)\colon\left(x-1\right)^2+\left(y+2\right)^2+z^2=16$. Xét tính đúng sai của các khẳng định sau:
	 \choiceTF
		{\True Điểm $A\left(1;2;-1\right)$ nằm bên ngoài mặt cầu $(C)$} 
		{\True Điểm $B\left(0;0;1\right)$ nằm  bên trong mặt cầu $(C)$}
		{Điểm $C\left(0;2;1\right)$ nằm  trên mặt cầu $(C)$}
		{\True Với điểm $D\left(2;1;-1\right)$,  ta có $ID<4$}
	\loigiai{ 
		\begin{itemchoice}
			\itemch Tâm $I\left(1;-2;0\right)$.\\ 
			Ta có $\left(1-1\right)^2+\left(2+2\right)^2+\left(-1\right)^2=17>16$ $\Rightarrow$ điểm $A\left(1;2;-1\right)$ nằm  bên ngoài mặt cầu $(C)$.
			 \itemch Ta có $\left(0-1\right)^2+\left(0+2\right)^2+1^2=6<16$ $\Rightarrow $ điểm $B\left(0;0;1\right)$ nằm bên trong mặt cầu $(C)$.
			\itemch Ta có $\left(0-1\right)^2+\left(2+2\right)^2+1^2=18>16$ $\Rightarrow$ điểm $C\left(0;2;1\right)$ nằm bên ngoài mặt  cầu $(C)$.
			 \itemch Ta có $ID=\sqrt{\left(2-1\right)^2+\left(1+2\right)^2+\left(-1\right)^2}=\sqrt{11}<4$.
		\end{itemchoice}
	 }
\end{ex}

% nguồn: Phan1_Lop_12_Hinh_hoc_Chuong_5
\begin{ex}%[2H5V3-2]%[ma:950ac148]
	Trong không gian $Oxyz$, cho mặt cầu  $(S)\colon (x-2)^2+(y+6)^2+z^2=50$ có tâm $I$ và điểm $K(1;-3; 0)$.
	\choiceTF 
		{\True Mặt cầu  $(S)$ có bán kính $R=5 \sqrt{2}$}
		{Điểm $K$ thuộc mặt cầu $(S)$} 
		{Khoảng cách  từ tâm $I$ của mặt cầu $(S)$ đến điểm $K$ bằng $10$}
		{\True Xét  các điểm $M$ thuộc $(S)$ sao cho $\widehat{K M I}$ lớn nhất. Khi đó, giá trị của $M K$ bằng $2 \sqrt{10}$}
	\loigiai{ 
		\begin{itemchoice}
			\itemch  \textbf{Đúng}.
			Bán kính  $R=\sqrt{50}=5 \sqrt{2}$.
			\itemch  \textbf{Sai}.
			Thay toạ  độ của điểm $K$ vào phương trình mặt cầu $(S)$, ta có $$(1-2)^2+(-3+6)^2+0^2=10 \neq 50.$$
			Do đó, điểm  $K$ không thuộc mặt cầu $(S)$.
			\itemch \textbf{Sai}. 
			Ta có $\overrightarrow{I K}=(1;-3; 0)  \Rightarrow I K=\sqrt{1^2+(-3)^2+0^2}=\sqrt{10}$.
			 \itemch \textbf{Đúng}.
			Mặt cầu  $(S)$ có tâm $I(2;-6; 0)$, bán kính $R=5 \sqrt{2}$.\\
			Vì $M$ thuộc $(S)$  nên $M I=R=5 \sqrt{2}$.\\
			Ta có $\cos \widehat{K M I}=\dfrac{M K^2+M I^2-K I^2}{2 \cdot M K \cdot M I}=\dfrac{M  K^2+50-10}{2 \cdot M K \cdot 5 \sqrt{2}}=\dfrac{M K^2+40}{10 \sqrt{2} \cdot M K}$.\\
			Vì $0^{\circ}<\widehat{K M I}<90^{\circ}$ nên $0<\cos  \widehat{K M I}<1$.\\
			 Đặt $x=M K$, $x>0$. Xét hàm số $f(x) =\dfrac{x^2+40}{10 \sqrt{2}x}$ với $x>0$.\\
			Hàm số $f(x)$ đạt giá trị nhỏ nhất trên  $(0;+\infty)$ bằng $\dfrac{2}{\sqrt{5}}$ tại $x=2 \sqrt{10}$.\\
			Do đó $\cos \widehat{K M I}$ đạt giá trị  nhỏ nhất khi $M K=2 \sqrt{10}$.\\
			Vậy $\widehat{K M I}$ lớn  nhất khi giá trị của $M K$ bằng $2 \sqrt{10}$.
		\end{itemchoice}
	 }
\end{ex}

% nguồn: Phan1_Lop_12_Hinh_hoc_Chuong_5
\begin{ex}%[2H5V3-2]%[ma:cefc2532]
	Trong không gian $Oxyz$, cho ba điểm $A(1; 2;  3)$, $B(-2; 1; 1)$ và $C(-3; 3;-2)$.
	 \choiceTF
		{Toạ độ của $\overrightarrow{AB}$ là $(3; 1; 2)$} 
		{\True  Độ dài của đoạn thẳng $AC$ bằng $\sqrt{42}$}
		{$\widehat{BAC}$  là góc tù}
		{\True Xét các điểm $M$ trên mặt phẳng $(Ozx)$ thỏa mãn điều kiện $\widehat{AMB}=90^{\circ}$. Khi đó, giá trị lớn nhất của độ dài đoạn thẳng $OM$ lớn  hơn hoặc bằng $3$}
	\loigiai{ 
		\begin{itemchoice}
			 \itemch Ta có $ \overrightarrow{AB}=(-2-1; 1-2; 1-3)=(-3;-1; -2)$.
			\itemch Ta có  $AC=\sqrt{(-3-1)^2+(3-2)^2+(-2-3)^2}=\sqrt{42}$.
			\itemch Ta có $\overrightarrow{AB}=(-3;  -1; -2)$, $\overrightarrow{AC}=(-4; 1; -5)$.\\
			Do  đó $\cos \widehat{BAC}=\dfrac{\overrightarrow{AB} \cdot \overrightarrow{AC}}{AB \cdot AC}=\dfrac{\sqrt{3}}{2} \Rightarrow \widehat{BAC}=30^{\circ}$.\\
			Vậy  $\widehat{BAC}$ là góc nhọn.
			 \itemch
			Do $\widehat{AMB}=90^{\circ}$ nên $M$  thuộc mặt cầu đường kính $A B$ có tâm $I\left(-\dfrac{1}{2}; \dfrac{3}{2}; 2\right)$ và bán kính $R=\dfrac{A B}{2}=\dfrac{\sqrt{14}}{2}$.\\
			Hơn nữa, $M$ thuộc mặt phẳng $(Ozx)$ nên các điểm $M$ thuộc đường tròn $(C)$ là giao của mặt cầu $(I; R)$ với mặt phẳng $(O z  x)$.\\
			Đường tròn $(C)$ có tâm $J\left(-\dfrac{1}{2}; 0; 2\right)$ và  bán kính $r$ là
			$$r=\sqrt{R^2-I  J^2}=\sqrt{\dfrac{7}{2}-\dfrac{9}{4}}=\dfrac{\sqrt{5}}{2}.$$
			Vậy giá trị lớn nhất của $OM$  là $OJ+r=\dfrac{\sqrt{17}}{2}+\dfrac{\sqrt{5}}{2} \approx 3{,}18$.
		\end{itemchoice}
	 }
\end{ex}

% nguồn: Phan1_Lop_12_Hinh_hoc_Chuong_5
\begin{ex}%[2H5V3-2]%[ma:8360a68a]
	Trong không gian $Oxyz$, cho  mặt cầu $(S)\colon\left(x+1\right)^2+\left(y-1\right)^2+\left(z+2\right)^2=9$ có tâm $I$ và bán kính $R$.
	 \choiceTF
		{Phương trình mặt cầu $(S)$ có tâm $I\left(1;-1;2\right)$,  bán kính $R=3$}
		{Điểm $A\left(0;2;-3\right)$ nằm trong mặt  cầu}
		 {\True Điểm $J\left(1;2;3\right)$ nằm ngoài mặt cầu và khoảng cách từ tâm $I$ đến điểm $J$ bằng $\sqrt{10}$}
		{Khoảng cách từ tâm $I$ đến  tâm mặt cầu $\left(S'\right)\colon{x^2}+y^2+\left(z-1\right)^2=3$ bằng $\sqrt{2}$}
	\loigiai{ 
		\begin{itemchoice}
			\itemch Phương trình mặt cầu tâm $I\left(1;-1;2\right)$, bán kính $R=3$ là $\left(x-1\right)^2+\left(y+1\right)^2+\left(z-2\right)^2=9$. 
			\itemch Ta có $\left(0+1\right)^2+\left(2-1\right)^2+\left(-3+2\right)^2=3=R=3$, do đó điểm $A$  nằm trên mặt cầu.
			\itemch Ta có $\left(1+1\right)^2+\left(2-1\right)^2+\left(3+2\right)^2=30>R=3$, do đó điểm $J$  nằm ngoài mặt cầu.\\
			Và  $\overrightarrow{IJ}=\left(0;3;1\right)\Rightarrow\left|\overrightarrow{IJ}\right|=\sqrt{0^2+3^2+1^2}=\sqrt{10}$.
			\itemch Mặt cầu $(S)$ có  tâm $I\left(1;-1;2\right)$.\\
			Mặt cầu $\left(S'\right)$ có  tâm $K\left(0;0;1\right)$.\\
			Khi đó  $\overrightarrow{IK}=\left(-1;1;-1\right)\Rightarrow\left|\overrightarrow{IK}\right|=\sqrt{\left(-1\right)^2+1^2+\left(-1\right)^2}=\sqrt{3}\ne\sqrt{2}$.
		\end{itemchoice}
	} 
\end{ex}

% nguồn: Phan1_Lop_12_Hinh_hoc_Chuong_5
\begin{ex}%[2H5V3-2]%[ma:5bdae678]
	Trong không gian $Oxyz$, cho ba điểm $A(0; 1; 1)$, $B(3; 0; -1)$, $C(0; 21; -19)$ và mặt  cầu $(S)\colon (x-1)^2+(y-1)^2+(z-1)^2=1$.
	 \choiceTF
		{Tâm $I$ của mặt cầu $(S)$ có tọa độ là $(-1; -1;  -1)$}
		 {\True Toạ độ của điểm $D$ sao cho $3\overrightarrow{DA}+2 \overrightarrow{DB}+\overrightarrow{DC}=\overrightarrow{0}$ là $(1; 4; -3)$}
		{\True Đường  thẳng $\Delta$ đi qua hai điểm $I$ và $D$ có một vectơ chỉ phương là $(0; 3; -4)$}
		{Gọi $M$ là điểm thuộc mặt cầu $(S)$  sao cho biểu thức $T=3MA^2+2MB^2+MC^2$ đạt giá trị nhỏ nhất. Khi đó, giá trị của $T$ bằng $672$}
	 \loigiai{
		\begin{itemchoice}
			\itemch Tâm $I$  của mặt cầu $(S)$ có toạ độ là $(1; 1; 1)$.
			 \itemch
			Gọi $D(x; y; z)$, ta có $3 \overrightarrow{DA}+2 \overrightarrow{DB}+\overrightarrow{DC}=(-6x+6; -6y+24;  -6z-18)$.\\
			Vì  $3\overrightarrow{DA}+2 \overrightarrow{DB}+\overrightarrow{DC}=\overrightarrow{0}\Leftrightarrow \heva{& -6x+6=0 \\ & -6y+24=0 \\& -6z-18=0}\Leftrightarrow \heva{& x=1 \\ & y=4\\&z=-3.}$\\
			Vậy $D(1; 4;  -3)$.
			 \itemch
			Đường thẳng $\Delta$ đi qua hai điểm $I$ và $D$ có một vectơ chỉ phương là $\overrightarrow{ID}=(0; 3; -4)$. 
			\itemch 
			\begin{eqnarray*}
				 T&=&3 M A^2+2 M B^2+M C^2\\
				&=&3 \overrightarrow{M  A}^2+2 \overrightarrow{M B}^2+\overrightarrow{M C}^2 \\
				& =&3(\overrightarrow{M D}+\overrightarrow{D  A})^2+2(\overrightarrow{M D}+\overrightarrow{D B})^2+(\overrightarrow{M D}+\overrightarrow{D C})^2 \\
				& =&6 \overrightarrow{M D}^2+2 \overrightarrow{M D}(3 \overrightarrow{D  A}+2 \overrightarrow{D B}+\overrightarrow{D C})+3 D A^2+2 D B^2+D C^2 \\
				& =&6 M D^2+3 D A^2+2 D B^2+D C^2  .
			\end{eqnarray*}
			Ta có $ID =5 > R=1$ suy ra điểm  $D(1; 4; -3)$ nằm ngoài mặt cầu $(S)$ nên biểu thức $T$ đạt giá trị nhỏ nhất khi $MD$ ngắn nhất.\\
			Ta có $MD$ ngắn nhất khi và  chỉ khi $MD=ID-R=5-1=4$.\\
			 Giá trị nhỏ nhất của biểu thức $T$ là
			 $$T=6\cdot4^2+3\cdot26+2\cdot24+546=768.$$
		\end{itemchoice}
	} 
\end{ex}


%%%%%%%%%%%%%%%% TRẢ LỜI NGẮN %%%%%%%%%%%%%%%%

% nguồn: Phan2_Lop_12_Hinh_hoc_Chuong_5
\begin{bt}%[2H5N3-2]%[ma:4aed5fbe]
	Trong không gian $Oxyz$, cho mặt cầu $(S) \colon x^2+(y-2)^2+(z+1)^2=6$. Đường kính của $(S)$ bằng bao nhiêu? (kết  quả làm tròn đến hàng phần mười).
	\shortans{4,9} 
	\loigiai{Bán kính của  mặt cầu $(S)$ là $R=\sqrt{6}$.\\
	Do đó đường kính của $(S)$  bằng $2\sqrt{6} \approx 4{,}9$.}
\end{bt}

% nguồn: Phan2_Lop_12_Hinh_hoc_Chuong_5
\begin{bt}%[2H5N3-2]%[ma:06f363fd]
	Trong không gian với hệ  toạ độ $Oxyz$, cho mặt cầu $( S)\colon {x^2}+{y^2}+{z^2}-2x-4y+6z-67=0$. Bán kính của mặt cầu bằng bao nhiêu?
	\shortans{9} 
	 \loigiai{
	Mặt cầu $(S)$ có tâm $I( 1;2;-3)$.  \\Vậy bán kính của mặt cầu là $R=\sqrt {1^2+2^2+(-3)^2+67}=9$.}
\end{bt}

% nguồn: Phan2_Lop_12_Hinh_hoc_Chuong_5
\begin{bt}%[2H5H3-2]%[ma:ae1f41b8]
	Trong không gian với hệ  tọa độ $Oxyz$, cho mặt cầu $(S)\colon x^2+y^2+z^2-8x+16y+21z-100=0$. Giả sử $I\left(a,b,c\right)$ là tâm của $(S)$. Giá trị của biểu thức $M=a-2b+c$ là bao nhiêu?
	 \shortans{9,5}
	\loigiai{Ta có nếu  phương trình mặt cầu có dạng $x^2+y^2+z^2-2ax-2by-2cz+d$ thì mặt cầu đó có tâm là $I\left(a,b,c\right)$.\\
		Do đó $a=4$, $b=-8$, $c=-\dfrac{21}{2}$  nên $I\left(4;-8;-\dfrac{21}{2}\right)$.\\
	Suy  ra $M=a-2b+c=4+16-\dfrac{21}{2}=9{,}5$.}
\end{bt}

% nguồn: Phan2_Lop_12_Hinh_hoc_Chuong_5
\begin{bt}%[2H5H3-2]%[ma:b2dc9a26]
	Ericsson Globe (Thụy Điển) là tòa nhà bán cầu lớn  nhất trên thế giới (năm 2020), với hình dạng một quả cầu màu trắng có đường kính là $110$ m và chiều cao bên trong $85$ m, nó có đủ chỗ ngồi cho $16\,000$ khán giả của các buổi biểu diễn hòa nhạc hoặc $13\,850$ khán giả của các trận đấu khúc côn cầu trên băng. Giả sử ta biểu diễn mô phỏng của tòa nhà Ericsson Globe trong hệ trục tọa độ $Oxyz$ bởi một mặt cầu có tâm $I$, đường kính $110$ m và $OA=85$ m như hình vẽ (đơn vị trên trục là mét). Phương trình của mặt cầu này có dạng $x^2+y^2+\left(z-a\right)^2=b^2$. Tính $a+b$.
	\begin{center}
		\begin{tikzpicture}[line join=round,line cap=round,>=stealth,font=\footnotesize,scale=1]
				 \def\a{6.5}
				 \def\b{2.3}
				\def\h{1.7} 
				\path (0:0) coordinate (E) 
					 ++(0:\a) coordinate (D)
					 ++(-135:\b) coordinate (C)
					 ($(E)+(C)-(D)$) coordinate (B)
					($(E)!1/2!(C)$) coordinate (O) 
					($(O)+(90:\h*0.7)$)  coordinate (I)
					($(I)+(90:\h)$) coordinate (A) 
					 ($(I)+(-45:\h)$) coordinate (F)
					($(I)+(-135:\h)$)  coordinate (G)
					 ($(O)+(-135:0.5)$) coordinate (H)
					 ($(E)!0.1!(D)$) coordinate (D1)
					 ($(E)!0.65!(D)$) coordinate (D2);
				\draw[dashed] (O)--(A) 
					 (D1)--(D2);
				\draw  (E)--(B)--(C)--(D)
					 (E)--(D1)
					 (D2)--(D);
				 \draw (A) arc (90:-45:\h);
				\draw (A)  arc (90:225:\h);
				\draw[dashed]  (F) arc
					(0:180:{0.7141*\h} and  {0.7141*\h*0.25});
				 \draw (G) arc
					(180:360:{0.7141*\h} and {0.7141*\h*0.25}); 
				\draw[dashed] (O)--(H) 
					 (O)--(F);
				\foreach \x/\g in  {O/180,I/30}
				\fill (\x)  circle (1pt)
					($(\g:3mm)+(\x)$)  node {$\x$};
				 \draw (A) node[above right]{$A$};
				\draw[->](H)--($(O)!3.5!(H)$) node [pos=0.9,below  left]{$x$};
				\draw[->](F)--($(O)!1.9!(F)$)  node [pos=0.9,above]{$y$};
				 \draw[->](A)--($(O)!1.3!(A)$) node [pos=0.9,above right]{$z$};
				\draw ($(B)!0.1!(C)$)  arc
					(0:100:0.3) node  [pos=0.2,above right] {Mặt đất};
			\end{tikzpicture}
	\end{center}
	 \shortans{85}
	\loigiai{Vì $I \in Oz$  nên $I(0;0;z_I)$ với $z_I=OI$.\\
		Vì mặt cầu  có đường kính $110$ m nên $IA=\dfrac{110}{2}=55$ m.\\
		Do đó $OI=OA-IA=85-55=30$ m  nên $I(0;0;30)$.\\
		 Phương trình mặt cầu có tâm $I(0;0;30)$ và bán kính $IA=55$ có dạng
		 $$x^2+y^2+(z-30)^2=55^2.$$
	Vậy $a=30$,  $b=55$ nên $a+b=85$.}
\end{bt}

% nguồn: Phan2_Lop_12_Hinh_hoc_Chuong_5
\begin{bt}%[2H5H3-2]%[ma:28f81e6a]
	Trong không gian với  hệ tọa độ $Oxyz$, cho mặt cầu $(S)$ có tâm $I\left(2;-4;3\right)$ và tiếp xúc với trục $Ox$. Bán kính của mặt cầu trên bằng bao nhiêu?
	 \shortans{5}
	 \loigiai{
		Gọi $R$ là bán kính của mặt cầu $(S)$ và $H$ là hình chiếu của $I$ trên trục $Ox$.\\ 
		Suy ra  $H\left(2;0;0\right)$ và $R=IH=\sqrt{\left(2-2\right)^2+\left(0+4\right)^2+\left(0-3\right)^2}=5$.\\
		Mặt cầu $(S)$  có tâm $I\left(2;-4;3\right)$ và bán kính là $R=5$.
	 }
\end{bt}

% nguồn: Phan2_Lop_12_Hinh_hoc_Chuong_5
\begin{bt}%[2H5H3-2]%[ma:efe99dd1]
	Trong không gian với hệ tọa độ $Oxyz$, gọi $(S)$ là mặt cầu đi qua hai điểm $A(3;-1;2)$, $B(1;1;-2)$ và có tâm thuộc trục $Oz$. Bán kính của mặt  cầu $(S)$ là $R$. Giá trị $R^2$ bằng?
	\shortans{11} 
	 \loigiai{
		 Gọi tâm của mặt cầu là $I(a;b;c)$. Vì $I\in Oz$ nên $I(0;0;c)$.\\
		 Mặt cầu đi qua $A$, $B$ nên có $IA=IB\Leftrightarrow IA^2=IB^2$ $\Leftrightarrow 9+1+\left(c-2\right)^2=1+1+\left(c+2\right)^2$ $\Leftrightarrow c=1$.\\
		 Bán kính mặt cầu $R=IA=\sqrt{11}$.\\
		Vậy giá trị $R^2=11$. 
	 }
\end{bt}

% nguồn: Phan2_Lop_12_Hinh_hoc_Chuong_5
\begin{bt}%[2H5H3-2]%[ma:047aacbe]
	Trong không gian với hệ tọa độ $Oxyz$, cho điểm $M\left(1;-2;3\right)$ và $I$ là hình chiếu vuông góc của $M\left(3;1;2\right)$ trên trục $Ox$. Gọi $(S)$ là mặt cầu có tâm $I$ và đi qua điểm $M$. Độ dài bán kính  mặt cầu bằng bao nhiêu?
	\shortans{3} 
	\loigiai{ 
	 Hình chiếu vuông góc của $M$ trên trục $Ox$ là $I\left(1;0;0\right)$ $\Rightarrow IM=\sqrt{2^2+1^2+2^2}=3$.}
\end{bt}

% nguồn: Phan2_Lop_12_Hinh_hoc_Chuong_5
\begin{bt}%[2H5H3-2]%[ma:df828971]
	Trong không gian $Oxyz$, gọi $(S)$ là mặt cầu có tâm $I$ thuộc trục hoành $Ox$ và đi qua hai điểm $A(-6;6;0)$, $B(-4;0;8)$. Tính bán  kính $R$ của mặt cầu $(S)$.
	\shortans{10} 
	\loigiai{ 
		Gọi tọa độ điểm  $I(a;0;0)$.\\
		Mặt cầu tâm $I$ đi qua hai điểm $A$ và $B$ nên \begin{center}
		$IA^2=IB^2\Leftrightarrow  (-6-a)^2+6^2=(-4-a)^2+8^2\Leftrightarrow 4a=8\Leftrightarrow a=2$.
	\end{center}
	Suy ra  $R=IA=\sqrt{(-6-2)^2+6^2}=10$.
} 
\end{bt}

% nguồn: Phan2_Lop_12_Hinh_hoc_Chuong_5
\begin{bt}%[2H5V3-2]%[ma:da899ba4]
	Cho hai điểm $A$, $B$ cố định trong không gian có độ dài $AB$  là $4$. Biết rằng tập hợp các điểm $M$ trong không gian sao cho $MA=3MB$ là một mặt cầu. Bán kính mặt cầu đó bằng?
	\shortans{1,5} 
	\loigiai{ 
		Ta có 
		\begin{eqnarray*}
			&&MA=3MB\\ 
			&\Leftrightarrow&{\left(\overrightarrow{MA}\right)}^2=9{\left(\overrightarrow{MB}\right)}^2\\ 
			 &\Leftrightarrow&{\left(\overrightarrow{MI}+\overrightarrow{IA}\right)^2}=9\left(\overrightarrow{MI}+\overrightarrow{IB}\right)^2\\
			&\Leftrightarrow& IA^2-9IB^2+2\overrightarrow{MI}\left(\overrightarrow{IA}-9\overrightarrow{IB}\right)=8MI^2.\quad  (1)
		\end{eqnarray*}
		Gọi $I$ thỏa mãn $\overrightarrow{IA}-9\overrightarrow{IB}=\overrightarrow{0}\Leftrightarrow  \overrightarrow{BI}=\dfrac{1}{8}\overrightarrow{AB}$ nên $IB=\dfrac{1}{2}$; $IA=\dfrac{9}{2}$.\\
		Từ $(1)$ suy ra $\Leftrightarrow 8MI^2=18\Leftrightarrow MI=\dfrac{3}{2}$.\\ 
		Suy ra $M$ thuộc mặt cầu  $(S)$ có tâm là $I$ và bán kính $R=\dfrac{3}{2}=1{,}5$.
	 }
\end{bt}

% nguồn: Phan2_Lop_12_Hinh_hoc_Chuong_5
\begin{bt}%[2H5V3-2]%[ma:180d29c7]
	Cho mặt cầu $(S)\colon (x-1)^2+(y+2)^2+(z-3)^2=25$ và đường thẳng $d\colon\dfrac{x-2}2=\dfrac{y+4}{1}=\dfrac{z-3}{3}$. Đường thẳng $d$ cắt mặt cầu  $(S)$ tại hai điểm phân biệt trong đó $M\left( x_M;y_M;z_M \right)$ với $x_M>0$. Giá trị biểu thức $T=x_M^2+y_M^2+z_M^2$ bằng bao nhiêu (làm tròn kết quả đến hàng phần chục)?
	 \shortans{70,5}
	 \loigiai{
		Thay tọa độ tâm $I(1;-2;3)$ vào phương trình đường thẳng $d$, ta được $-\dfrac{1}2\ne  \dfrac2{1}\ne 0$.\\
		Nên  đường thẳng $d$ không đi qua tâm $I(1;-2;3)$.\\
		Ta có  đường thẳng $d\colon\heva{&x=2+2t\\& y=-4+t\\& z=3+3t}$ ($t\in \mathbb{R}$) với $M\in d\Rightarrow M(2+2t;-4+t;3+3t)$.\\
		Điểm  $M\in (S)$ nên ta có
		 \allowdisplaybreaks
		\begin{eqnarray*}
			 &&(2+2t-1)^2+(-4+t+2)^2+(3+3t-3)^2=25\\
			&\Leftrightarrow& (2t+1)^2+(t-2)^2+(3t)^2=25  \\
			&\Leftrightarrow& 14t^2+5=25\\ 
			&\Leftrightarrow& t^2=\dfrac{10}{7}\\ 
			&\Leftrightarrow&  \hoac{
			&t=\dfrac{\sqrt{70}}{7}\Rightarrow M\left( 2+\dfrac{2\sqrt{70}}{7};-4+\dfrac{\sqrt{70}}{7};3+\dfrac{3\sqrt{70}}{7}  \right) \\
			& t=-\dfrac{\sqrt{70}}{7}\Rightarrow N\left(  2-\dfrac{2\sqrt{70}}{7};-4-\dfrac{\sqrt{70}}{7};3-\dfrac{3\sqrt{70}}{7} \right)\text{ (loại vì }x=2-\dfrac{2\sqrt{70}}{7}<0).}
		\end{eqnarray*}
		Vậy $x_M^2+y_M^2+z_M^2=\left(2+\dfrac{2\sqrt{70}}{7} \right)^2+\left(-4+\dfrac{\sqrt{70}}{7} \right)^2+\left( 3+\dfrac{3\sqrt{70}}{7} \right)^2\approx  70{,}5$.
	 }
\end{bt}

% nguồn: Phan2_Lop_12_Hinh_hoc_Chuong_5
\begin{bt}%[2H5V3-2]%[ma:f81cde93]
	 Trong không gian $Oxyz$, một vòm được thiết kế có bề mặt là mặt cầu tâm $I(1;2;20)$, bán kính bằng $50$ m và có đáy nằm trên mặt phẳng $(Oxy)$. Chiều cao của vòm là bao nhiêu? (biết đơn vị của hệ trục tọa độ là mét).
	 \shortans{70}
	\loigiai{ 
		Khoảng cách từ tâm  $I$ đến mặt phẳng $(Oxy)$ là $20$ m.\\
		Vậy chiều  cao của vòm là $20+50=70$ m.
	} 
\end{bt}

% nguồn: Phan2_Lop_12_Hinh_hoc_Chuong_5
\begin{bt}%[2H5V3-2]%[ma:d0bf7ec7]
	Trong không gian $Oxyz$ cho mặt cầu $(S)\colon (x-1)^2+(y-2)^2+(z+1)^2=6$ và hai mặt phẳng $(P)\colon x+y+2z+5=0$, $(Q)\colon 2x-y+z-5=0$. Hai đường thẳng $d$ và $\Delta$ đi qua tâm của $(S)$, lần lượt vuông góc và  cắt $(P)$, $(Q)$ theo thứ tự đó tại $A$ và $B$. Tính độ dài đoạn thẳng $AB$ (kết quả làm tròn đến hàng phần trăm).
	\shortans{4,24} 
	 \loigiai{
		Mặt  cầu $(S)$ có tâm $I(1;2;-1)$. \\
		Phương trình đường thẳng $d$ và $\Delta$ lần lượt là $d\colon \dfrac{x-1}{1}=\dfrac{y-2}{1}=\dfrac{z+1}{2}$,  \break $\Delta\colon\dfrac{x-1}{2}=\dfrac{y-2}{-1}=\dfrac{z+1}{1}$. \\
		Ta có $d\cap(P)=A(0;1;-3)$,  $d\cap(Q)=B(3;1;0)$.\\
		Suy ra $\overrightarrow{AB}=(3;0;3)$. \\ 
		Do đó $AB=\left|\overrightarrow{AB}\right|=\sqrt{3^2+0^2+3^2}=3\sqrt{2}\approx 4{,}24$. 
	} 
\end{bt}

% nguồn: Phan2_Lop_12_Hinh_hoc_Chuong_5
\begin{bt}%[2H5V3-2]%[ma:d425100b]
	Trong không gian với hệ trục $Oxyz$, cho mặt phẳng $(P)\colon x + 2y - z - 3 = 0$, mặt cầu $(S)$ có tâm $I$ thuộc $(P)$ và tiếp xúc với $(Q)\colon x + y - z - 1 = 0$ tại $A(1; 1; 1)$, $(P)$ cắt mặt cầu $(S)$ theo đường tròn $(C)$. Tính bán kính của $(C)$ (kết quả làm tròn đến hàng  phần mười).
	\shortans{0,4} 
	\loigiai{ 
		\begin{itemize}
			 \item Tọa độ tâm $I$ của mặt cầu $(S)$.\\
			Vì $(S)$ tiếp xúc với $(Q)$ tại $A(1;1;1)$ nên $IA \perp (Q)$. Do đó,  vectơ pháp tuyến của $(Q)$ là $\overrightarrow{n_Q} = (1;1;-1)$ cũng là vectơ chỉ phương của $IA$.\\
			Phương trình đường thẳng $IA$ 
			 $$
			\begin{cases}
				x =  1 + t \\
				y  = 1 + t \\
				z  = 1 - t.
			\end{cases}
			 $$
			Vì $I \in (P)$ nên thay vào phương trình  $(P)$
			 $$
			(1 + t) + 2(1 + t) - (1 - t) -  3 = 0 \Leftrightarrow 4t - 1 = 0 \Leftrightarrow t = \dfrac{1}{4}.
			$$ 
			Suy ra  tọa độ tâm $ I\left(\dfrac{5}{4}; \dfrac{5}{4}; \dfrac{3}{4}\right)$.
			\item Bán  kính $R$ của mặt cầu $(S)$.\\
			Bán  kính $R = IA = \sqrt{\left(\dfrac{5}{4}-1\right)^2 + \left(\dfrac{5}{4}-1\right)^2 + \left(\dfrac{3}{4}-1\right)^2} = \dfrac{\sqrt{3}}{4}$.
			\item Vì $I\in (P)$ nên đường tròn $(C)$ có bán kính bằng bán kính  mặt cầu $(S)$ và bằng $\dfrac{\sqrt{3}}{4}\approx 0{,}4$.
		\end{itemize}
	 }
\end{bt}

% nguồn: Phan2_Lop_12_Hinh_hoc_Chuong_5
\begin{bt}%[2H5V3-2]%[ma:ac9ed746]
	Trong không gian $Oxyz$, cho mặt cầu $(S)$ tâm  $I(-2;3;0)$ và đường thẳng $\Delta\colon\dfrac{x-4}2=\dfrac{y-3}{1}=\dfrac{z-3}2$. Biết đường thẳng $\Delta$ cắt mặt cầu $(S)$ tại hai điểm phân biệt $A$, $B$ sao cho $AB=8$. Tính bán kính của mặt cầu $(S)$.
	 \shortans{5}
	\loigiai{ 
		Giả sử mặt cầu $(S)$ có  phương trình $(S)\colon x^2+y^2+z^2+4x-6y+2m=0$.\\
		Mặt cầu $(S)$ có tâm $I(-2;3;0)$  và bán kính $R=\sqrt{13-2m}$ với $m<\dfrac{13}2$.\\
		Đường thẳng $\Delta$ đi qua điểm $M(4;3;3)$ và có một vectơ chỉ phương là $\overrightarrow{u}=(2;1;2)$.\\ 
		 Ta có $\overrightarrow{IM}=(6;0;3)$; $\left[\overrightarrow{IM},\overrightarrow{u} \right]=(-3;-6;6)$.\\
		Khoảng cách từ tâm $I$ đến đường thẳng $\Delta$ bằng 
		$$d=\mathrm{d}(I,\Delta)=\dfrac{\left| \left[\overrightarrow{IM},\overrightarrow{u} \right] \right|}{\left| \overrightarrow{u}  \right|}=\dfrac{\sqrt{(-3)^2+(-6)^2+6^2}}{\sqrt{2^2+1^2+2^2}}=3.$$
		\begin{center}
			\begin{tikzpicture}[scale=1,font=\footnotesize,line join=round,line cap=round,>=stealth]
					 \def \r{2}
					\path 
						(0,0) coordinate (I) 
						(160:\r) coordinate  (A)
						(60:\r) coordinate (B) 
						($(A)!(I)!(B)$) coordinate (H) 
						(180:\r) coordinate  (N)
						(intersection  of A--B and I--N) coordinate (M) node[above]{$d$};
					\draw  (I) circle (\r);
					\draw  (I)--(H)
						(M)--(B) 
						 (A)--(I)--(B);
					\pic[draw,angle radius=2mm] {right angle = I--H--B}; 
					\foreach  \x/\g in {A/150,B/60,I/-90,H/90}\fill (\x) circle (1pt)+(\g:.3) node{$\x$};
				\end{tikzpicture}
		\end{center}
		Vì đường thẳng $\Delta$ cắt mặt cầu $(S)$ tại hai điểm phân biệt $A$, $B$ sao cho $AB=8$ nên  ta có $d^2+\dfrac{AB^2}{4}={R^2}\Leftrightarrow 3^2+16=13-2m\Leftrightarrow m=-6$. (thỏa điều kiện).\\
		Suy ra bán kính mặt  cầu là $R=\sqrt{13-2\cdot (-6)}=5$.
	 }
\end{bt}

% nguồn: Phan2_Lop_12_Hinh_hoc_Chuong_5
\begin{bt}%[2H5V3-2]%[ma:4282fd46]
	Mặt cầu $(S)$ có tâm $I\left(-1;2;3\right)$ và có đúng một điểm chung với mặt phẳng $\left(Oxy\right)$. Gọi một giao điểm của  mặt cầu $(S)$ với trục $Oz$ là $A\left(a;b;c\right)$ với $\left(c>1\right)$. Tổng $a+b+c$ bằng bao nhiêu?
	\shortans{5} 
	\loigiai{ 
		Ta có phương trình mặt phẳng  $\left(Oxy\right)$ là $z=0$.\\
		Mặt cầu $(S)$ có đúng một điểm chung với  mặt phẳng $\left(Oxy\right)$ $\Leftrightarrow R=\mathrm{\,d}\left(I,\left(Oxy\right)\right)=3$.\\
		 Suy ra $(S):\left(x+1\right)^2+\left(y-2\right)^2+\left(z-3\right)^2=9$.\\
		Cho $x=y=0\Rightarrow {1^2}+2^2+\left(z-3\right)^2=9\Leftrightarrow \hoac{&z=5  \\&z=1.}$\\
		 Vì $c>1$ nên $A\left(0;0;5\right)\Rightarrow a+b+c=5$.
	} 
\end{bt}

% nguồn: Phan2_Lop_12_Hinh_hoc_Chuong_5
\begin{bt}%[2H5V3-2]%[ma:7ed17010]
	Trong không gian với hệ tọa độ $Oxyz$, mặt cầu $(S)$ có bán kính bằng $2$, tiếp xúc với mặt phẳng $\left(Oyz\right)$ và có tâm nằm  trên tia $Ox$. Hoành độ của tâm mặt cầu $(S)$ bằng bao nhiêu?
	 \shortans{2}
	\loigiai{ 
		 Gọi $I$ là tâm mặt cầu $(S)$. Vì tâm $I$ nằm trên tia $Ox$ nên $I\left(a;0;0\right)$, $\left(a>0\right)$.\\
		Vì mặt cầu $(S)$  có bán kính bằng $2$, tiếp xúc với mặt phẳng $\left(Oyz\right)$ nên ta có
		\begin{center}
			$\mathrm{\,d}\left(I,\left(Oyz\right)\right)=2 
			\Leftrightarrow \dfrac{\left|  a\right|}{1}=2
			\Leftrightarrow \left|  a\right|=2
			 \Leftrightarrow \hoac{&a=2 \\&a=-2.}$
		\end{center}
		Do $a>0$ nên ta nhận $a=2$, loại $a=-2$ suy ra  mặt cầu $(S)$ có tâm $I\left(2,0,0\right)$.
	} 
\end{bt}

% nguồn: Phan2_Lop_12_Hinh_hoc_Chuong_5
\begin{bt}%[2H5V3-2]%[ma:8d64fc81]
	Trong không gian với hệ tọa độ Oxyz, cho mặt phẳng $(P)\colon x+y-z-3=0$ và  hai điểm $M\left(1;1;1\right)$, $N\left(-3;-3;-3\right)$. Mặt cầu $(S)$ đi qua $M$, $N$ và tiếp xúc với mặt phẳng $(P)$ tại điểm $Q$. Biết rằng $Q$ luôn thuộc một đường tròn cố định. Tìm bán kính của đường tròn đó.
	 \shortans{6}
	 \loigiai{
		\begin{center}
			\begin{tikzpicture}[scale=1]
					\fill[gray!20] (-6,-4) -- (-4,-2) -- (5,-2) -- (3,-4)  -- cycle;
					\node at (4.2,-2.3) {$(P)$}; 
					 \def\r{3}
					\path (-4,-3) coordinate (A) 
						(0,-3) coordinate (Q) 
						(0,0) coordinate  (I)
						 ($(I)!3/5!(A)$) coordinate (H)
						(-1.5,2.598)  coordinate (N)
						($(A)!2.12/5!(N)$)  coordinate (M);
					\draw 
						 (0:0) circle (\r);
					\draw[dashed] 
						(180:\r) arc  (180:0:{\r} and {0.3*\r});
					 \draw
						(180:\r) arc (180:360:{\r} and  {0.3*\r});
					 \draw (A)--(Q)
						(A)--(H) 
						 (A)--(M);
					\draw[dashed]  (I)--(Q)
						(I)--(H) 
						 (M)--(N);
					 \foreach \x/\g in {A/-90,Q/-90,I/90,N/100,M/180}
					 \fill (\x) circle (1pt)
						($(\g:3mm)+(\x)$) node {$\x$}; 
				\end{tikzpicture}
		\end{center}
		Đường  thẳng $MN$ đi qua điểm $M\left(1;1;1\right)$ và nhận $\overrightarrow{MN}=\left(-4;-4;-4\right)$ là véctơ chỉ phương.\\
		Chọn $\overrightarrow{u}=\left(1;1;1\right)$ là véctơ chỉ phương của đường thẳng $MN$, suy ra phương trình đường  thẳng $MN$ là $\heva{& x=1+t\\
		& y=1+t\\ 
		& z=1+t.}$\\ 
		Gọi $A=MN\cap(P)$, suy ra tọa độ  $A$ là nghiệm hệ phương trình $\heva{
		&  x=1+t\\
		& y=1+t\\ 
		 & z=1+t\\
		&  x+y-z-3=0}\Leftrightarrow\heva{
		&  t=2\\
		&  x=3\\
		& y=3\\ 
		&  z=3.}$\\
		 Suy ra $A\left(3;3;3\right)$.\\
		Các điểm $M$, $N$, $Q$ cùng thuộc một đường tròn và $AQ$ là tiếp tuyến đường tròn  đó nên ta có $AQ^2=AM\cdot AN$.\\
		Thay tọa độ các điểm vào  ta được $AM=2\sqrt{3}$, $AN=6\sqrt{3}$, suy ra $AQ^2=36\Leftrightarrow AQ=6$.\\
	Vậy điểm $Q$  luôn thuộc một đường tròn cố định có tâm$A\left(3;3;3\right)$ và bán kính $R=6$.}
\end{bt}

% nguồn: Phan2_Lop_12_Hinh_hoc_Chuong_5
\begin{bt}%[2H5V3-2]%[ma:f9650300]
	Trong không gian $Oxyz$, mặt cầu có tâm $I(1;2;-1)$ và  tiếp xúc với mặt phẳng $(P)\colon 2x-2y-z-8=0$ có phương trình $(C)\colon x^2+y^2+z^2-2ax-2by-2cz+d=0$. Tính giá trị của biểu thức $a+b+c+2d$.
	 \shortans{-4}
	\loigiai{ 
		Bán kính mặt cầu là  $R=\mathrm{d}\left(I,(P)\right)=\dfrac{\left| 2-4+1-8 \right|}{\sqrt{4+4+1}}=3$.\\
		Phương trình mặt  cầu $(C)\colon (x-1)^2+(y-2)^2+\left(z+1 \right)^2=9 \Leftrightarrow x^2+y^2+z^2-2x-4y+2z-3=0$.\\
		 Suy ra $a=1$; $b=2$; $c=-1$; $d=-3$.\\
		Hay $a+b+c+2d=1+2+(-1)+2\cdot(-3)=-4$. 
	 }
\end{bt}

% nguồn: Phan2_Lop_12_Hinh_hoc_Chuong_5
\begin{bt}%[2H5V3-2]%[ma:58216a44]
	Trong không gian $Oxyz$, cho mặt phẳng $(P)\colon x+y-z-3=0$ và hai điểm $M(1;1;1)$, $N(-3;-3;-3)$. Mặt cầu $(S)$ đi qua hai điểm $M$, $N$ và tiếp xúc với mặt phẳng $(P)$ tại điểm $Q$. Biết rằng $Q$ luôn thuộc một đường tròn cố định. Tìm  độ dài bán kính của đường tròn đó.
	 \shortans{6}
	 \loigiai{
		\begin{center}
			\begin{tikzpicture}[line join=round,line cap=round,>=stealth,font=\footnotesize,scale=0.8]
					 \def\r{2cm}
					 \path
						(0,0) coordinate  (K)
						(-90:\r)  coordinate (Q)
						(90:\r)  coordinate (M)
						(-10:\r) coordinate  (N)
						($(M)!2!(N)$) coordinate  (x)
						 (Q)++(0:2*\r) coordinate (y)
						(intersection  of M--x and Q--y) coordinate (I)
						 (-30:\r) coordinate (E)
						(-150:\r)  coordinate (F)
						 ($(E)!1.1!(F)$) coordinate (B)
						($(E)!-1.2!(F)$)  coordinate (A)
						 (B)++(-130:\r) coordinate (C)
						 ($(A)+(C)-(B)$) coordinate (D)
						 (intersection of A--B and M--I) coordinate (G);
					 \draw[name path=c] (K) circle (\r);
					\path[name path=d]  (I)--(K);
					\path [name intersections={of=d and  c,by=X}];
					\draw[dashed] (M).. controls+ (-10:0.7) and +(90:0.7)..  (N)
						 (M)--(N)
						 (Q)--(K)--(X)
						 (F)--(G);
					 \draw
						(M)..  controls+ (-170:0.8) and +(-100:0.8).. (N)
						(N)--(I)--(Q) 
						(I)--(X) 
						(G)--(A)--(D)--(C)--(B)--(F); 
					 \node at ($(C)+(20:0.7)$) {$(P)$};
					\pic[draw,angle radius=0.2cm]{right  angle=I--Q--K};
					\foreach \i/\g in  {M/90,N/0,K/90,Q/-90,I/90}{\draw (\i) circle (1pt)
						 ($(\i)+(\g:3mm)$) node[scale=1]{$\i$};}
				\end{tikzpicture}
		\end{center}
		Đường thẳng $MN$ có phương trình $MN\colon\heva{&x=1+t\\& y=1+t\\&  z=1+t.}$\\
		Gọi $I=MN\cap(P)$, khi đó điểm $I$  ứng với $t$ thỏa mãn $$1+t+1+t-1-t-3=0\Leftrightarrow t-2=0\Leftrightarrow t=2.$$
		$\Rightarrow  I(3;3;3)\Rightarrow \heva{&IM=2\sqrt{3}\\&IN=6\sqrt{3}.}$\\
		Do mặt cầu $(S)$ đi qua $M$, $N$ và tiếp xúc với đường thẳng $IQ$ tại điểm $Q$  nên ta có
		\begin{align*}
			IQ^2&=IM\cdot  IN=KI^2-R^2\\
			 \Rightarrow IQ^2&=IM\cdot IN=36\\
			\Leftrightarrow  IQ&=6
		\end{align*}
		Vậy $Q$ luôn thuộc đường tròn tâm $I$ bán kính  $R=6$.
	} 
\end{bt}

% nguồn: Phan2_Lop_12_Hinh_hoc_Chuong_5
\begin{bt}%[2H5V3-2]%[ma:d3cb97b4]
	Trong không gian $Oxyz$, cho đường thẳng $d\colon \dfrac{x+1}{3}=\dfrac{y-2}{-2}=\dfrac{z-2}{2}$. Tìm bán kính mặt cầu $(S)$ có tâm $I(1; 2;-1)$ cắt $d$ tại các điểm $A$,  $B$ sao cho $AB=2\sqrt{3}$.
	\shortans{4} 
	\loigiai{ 
	\immini{ 
		Gọi $H$ là trung điểm của  $AB$, suy ra $HA=\sqrt{3}$.\\
		Vì $A$, $B$ thuộc mặt cầu $(S)$ tâm $I$ nên $IA=IB$, suy ra tam giác $IAB$ cân tại $I$, có $IH$ là đường trung tuyến, suy ra $IH$ cũng là  đường cao, suy ra $IH\perp AB$ tại $H\Rightarrow \vec{IH}\cdot \vec{u}=0$ (với $\vec{u}$ là véc-tơ chỉ phương của $d$).\\
		Mà  $H\in(d)$ nên $H(-1+3t;2-2t;2+2t)\ \left(t\in\mathbb{R}\right)$.\\
		Suy ra $\vec{IH}=(3t-2;-2t;2t+3)$;  $\vec{u}=(3;-2;2)$. Do đó
		\allowdisplaybreaks 
		\begin{eqnarray*}
			&&\vec{IH}\cdot \vec{u}=0\\ 
			&\Leftrightarrow &3(3t-2)-2\cdot  (-2t)+2(2t+3)=0\\
			&\Leftrightarrow &17t=0\\ 
			&\Leftrightarrow  &t=0.
		\end{eqnarray*}
		 }{
		\begin{tikzpicture}[>=stealth,line join=round,line cap=round,font=\footnotesize,scale=1]
				\def\r{3}; 
				\def\g{40}; 
				\pgfmathsetmacro\a{\r*sin  (90-\g)};
				\path 
					(0,0) coordinate  (I)
					(\g-180:\r) coordinate  (A)
					(-\g:\r) coordinate (B) 
					($(A)!(I)!(B)$)  coordinate (H);
				\pgfresetboundingbox 
				\draw[dashed] (0:\r) arc  (0:180:\r cm and 0.2*\r cm)
					(A)--(I)--(B)--cycle  (I)--(H);
				\draw (I) circle (\r) 
					(180:\r) arc  (180:360:\r cm and 0.2*\r cm);
				\foreach \x/\g in {A/-135,I/90,B/-45,H/-90}\fill (\x) circle (1pt)  +(\g:3mm) node {$\x$};
				\foreach  \x/\y in {A/H,H/B}{\node at ($(\x)!0.5!(\y)$){$\parallel$};}
			\end{tikzpicture}
		} 
		\noindent Với $t=0$ suy  ra $H=(-1;2;2)$, $\vec{IH}=(-2;0;3)\Rightarrow IH=\sqrt{4+0+9}=\sqrt{13}$.\\
		Áp dụng định lý Pythagore vào  tam giác vuông $IHA$, ta có
		 \[IA=\sqrt{IH^2+HA^2}=\sqrt{13+3}=4. \]
		Vậy  bán kính mặt cầu $(S)$ bằng $4$ (đơn vị độ dài).
	} 
\end{bt}

% nguồn: Phan2_Lop_12_Hinh_hoc_Chuong_5
\begin{bt}%[2H5V3-2]%[ma:2d6c4d40]
	Trong không gian với hệ trục tọa độ $Oxyz$, cho hai  điểm $A(1;0;-3)$, $B(-3;-2;-5)$. Biết rằng tập hợp các điểm $M$ trong không gian thỏa mãn đẳng thức $AM^2+BM^2=30$ là một mặt cầu $(S)$. Tìm bán kính $R$ của mặt cầu $(S)$.
	 \shortans{3}
	 \loigiai{
		Gọi tọa độ điểm $M(x;y;z)$.\\ 
		Khi  đó
		\begin{eqnarray*}
			&&AM^2+BM^2=30\\ 
			&\Leftrightarrow&{\left(x-1\right)^2}+y^2+\left(z+3\right)^2+\left(x+3\right)^2+\left(y+2\right)^2+\left(z+5\right)^2=30\\ 
			 &\Leftrightarrow& 2x^2+2y^2+2z^2+4x+4y+16z+18=0\\
			 &\Leftrightarrow& x^2+y^2+z^2+2x+2y+8z+9=0\\
			 &\Leftrightarrow&{\left(x+1\right)^2}+\left(y+1\right)^2+\left(z+4\right)^2=9.
		\end{eqnarray*}
		Nên mặt cầu $(S)$ có tâm $I\left(-1;-1;-4\right)$ và bán  kính $R=3$.
	 }
\end{bt}

% nguồn: Phan2_Lop_12_Hinh_hoc_Chuong_5
\begin{bt}%[2H5V3-2]%[ma:08bc9517]
	Bán kính mặt cầu đi qua bốn điểm $M(1;0;1)$, $N(1;0;0)$, $P(2;1;0)$ và $Q(1;1;1)$ có dạng $\dfrac{\sqrt{m}}{n}$; $m,\,n\in \mathbb{N}^*$ và nguyên  tố cùng nhau. Khi đó giá trị của biểu thức $S=m^n$ bằng bao nhiêu?
	\shortans{9} 
	\loigiai{ 
		Gọi  phương trình mặt cầu $(S)\colon x^2+y^2+z^2-2ax-2by-2cz+d=0$, ($a^2+b^2+c^2-d>0$).\\
		 Do $(S)$ đi qua bốn điểm $M$, $N$, $P$, $Q$ nên ta có hệ phương trình
		$$\heva{&  -2a-2c+d=-2\\& -2a+d=-1\\&-4a-2b+d=-5\\&-2a-2b-2c+d=-3}\Leftrightarrow \heva{& a=\dfrac{3}2\\& b=\dfrac{1}2\\& c=\dfrac{1}2\\& d=2.}$$
		Vậy $R=\sqrt{\left(\dfrac{3}2\right)^2+\left(\dfrac{1}2^2\right)+\left(\dfrac{1}2\right)^2-2}=\dfrac{\sqrt{3}}2$.\\ 
		Suy ra  $m=3$, $n=2$.\\
		Khi đó giá trị của biểu thức $S=m^n=3^2=9$. 
	 }
\end{bt}

% nguồn: Phan2_Lop_12_Hinh_hoc_Chuong_5
\begin{bt}%[2H5V3-2]%[ma:554e4caa]
	Trong không gian $Oxyz$, cho mặt cầu $(S)\colon x^2+y^2+z^2+4x-6y+m=0$ ($m$ là tham số) và đường thẳng $\Delta\colon\heva{&x=4+2t\\ &y=3+t\\ &z=3+2t}$. Biết đường thẳng $\Delta$ cắt mặt cầu $(S)$ tại hai điểm phân  biệt $A$, $B$ sao cho $AB=8$. Tính giá trị của $m$.
	 \shortans{-12}
	\loigiai{ 
		\begin{center}
			\begin{tikzpicture}[line join=round,line cap=round,>=stealth,font=\footnotesize,scale=1]
					\def\r{2cm} 
					 \path
						(0,0)  coordinate (I)
						(35:\r) coordinate  (B)
						(145:\r)  coordinate (A)
						($(A)!1.3!(B)$) coordinate (x) 
						 ($(A)!(I)!(B)$) coordinate (H)
						 ($(B)!1.3!(A)$) coordinate (y)
						 ($(B)!1.2!(A)$) coordinate (M);
					 \draw
						(y)--(x) node[above]{$\Delta$} 
						 (I) circle (\r)
						 (B)--(I) node[midway,sloped,above]{$R$}--(A)
						(I)--(H); 
					\pic[draw,angle  radius=0.2cm]{right angle=I--H--B};
					\foreach \i/\g  in {H/90,A/120,B/70,M/90,I/-90}{\draw (\i) circle (1pt)
						($(\i)+(\g:3mm)$)  node[scale=1]{$\i$};}
				\end{tikzpicture}
		\end{center}
		Gọi $H$ là trung điểm đoạn thẳng $AB\Rightarrow IH\perp  AB$, $HA=4$.\\
		 Mặt cầu $(S)$ có tâm $I(-2;3;0)$, bán kính $R=\sqrt{13-m},\,\,(m < 13)$.\\
		Đường thẳng $\Delta$ đi qua $M(4;3;3)$ và có vectơ  chỉ phương $\overrightarrow{u}=(2;1;2)$.\\
		Ta  có
		 $\overrightarrow{IM}=(6;0;3)\Rightarrow\left[\overrightarrow{IM},\overrightarrow{u}\right]=(-3;-6;6)$\\
		 $\Rightarrow IH=d_{(I,\Delta)}=\dfrac{\left|\left[\overrightarrow{IM},\overrightarrow{u}\right]\right|}{|\overrightarrow{u}|}=3$.\\
		Ta có: $R^2=IH^2+HA^2\Leftrightarrow 13-m=3^2+4^2\Leftrightarrow m=-12$. 
	} 
\end{bt}

% nguồn: Phan2_Lop_12_Hinh_hoc_Chuong_5
\begin{bt}%[2H5V3-2]%[ma:b9dd12fe]
	Trong không gian $Oxyz$, mặt cầu $(S)$ có tâm là điểm $A\left(1;2;-3\right)$ và đi qua điểm $B\left(3;-2;1\right)$. Gọi $M\left(1;b;-3\right)$ là điểm thuộc mặt cầu biết $b\in{\mathbb{N}^{*}}$. Tung độ  điểm $M$ có giá trị bằng bao nhiêu?
	\shortans{8} 
	 \loigiai{
		Mặt cầu $(S)$ đi  qua $B\left(3;-2;1\right)$ nên bán kính $R=AB$
		$\Leftrightarrow R=\sqrt{\left(3-1\right)^2+\left(-2-2\right)^2+\left(1+3\right)^2}\Leftrightarrow R=6$.\\ 
		Vậy phương trình của  mặt cầu $(S)$ là: $\left(x-1\right)^2+\left(y-2\right)^2+\left(z+3\right)^2=36$.\\
		Thay tọa độ  điểm $M\left(1;b;-3\right)$ vào phương trình mặt cầu $(S)$, ta có
		\begin{center}
			 $\left(1-1\right)^2+\left(b-2\right)^2+\left(-3+3\right)^2=36
			\Leftrightarrow{\left(b-2\right)^2}=36 
			\Leftrightarrow  \hoac{&b-2=6 \\&b-2=-6}
			\Leftrightarrow \hoac{&b=8  \\&b=-4.}$
		\end{center}
		 Vì $b>0$ nên nhận $b=8$.
	} 
\end{bt}

% nguồn: Phan2_Lop_12_Hinh_hoc_Chuong_5
\begin{bt}%[2H5V3-2]%[ma:cd6af2e4]
	Trong không gian $Oxyz$, phương trình mặt cầu $(S)$ có tâm $I(1;3;2)$ và có một tiếp tuyến là đường  thẳng $\Delta\colon \dfrac{x-1}{1}=\dfrac{y-1}{1}=\dfrac{z}{-1}$ có dạng $x^2+y^2+z^2-2ax-2by-2cz+d=0$. Tính $T=a+b+c-d$.
	 \shortans{0}
	 \loigiai{
		Ta có đường thẳng $\Delta\colon\heva{&M_0(1;1;0)\in \Delta\\& 
		\overrightarrow{u}=(1;1;-1)\text{ là vectơ chỉ  phương.}}$\\
		 Suy ra $\overrightarrow{IM}_0=(0;-2;-2)$.\\
		Ta có  $\left[\overrightarrow{u},\overrightarrow{IM}_0 \right]=(-4;2;-2)$.\\
		Vì $\Delta$ là một tiếp tuyến của mặt cầu $(S)$  nên
		 $$R=\mathrm{d}\left( M,\Delta \right)=\dfrac{\left| \left[\overrightarrow{u},\overrightarrow{IM}_0 \right] \right|}{\left| {\overrightarrow{u}} \right|}=\dfrac{\sqrt{(-4)^2+2^2+(-2)^2}}{\sqrt{1^2+1^2+(-1)^2}}=\dfrac{\sqrt{24}}{\sqrt{3}}=2\sqrt2.$$
		Khi đó phương  trình mặt cầu $(S)$ có dạng
		\allowdisplaybreaks 
		\begin{eqnarray*}
			&&(x-1)^2+(y-3)^2+(z-2)^2=8\\ 
			&\Leftrightarrow& x^2+y^2+z^2-2x-6y-4z+6=0\Leftrightarrow\heva{&a=1\\& b=3\\ 
			& c=2 \\&  d=6.}
		\end{eqnarray*}
		Suy ra $T=a+b+c-d=0$. 
	} 
\end{bt}

% nguồn: Phan2_Lop_12_Hinh_hoc_Chuong_5
\begin{bt}%[2H5V3-2]%[ma:a942b0fe]
	Trong không gian $Oxyz$, cho mặt cầu $(S)$ có tâm $I(a;b;c)$ nằm trên đường thẳng $d\colon\dfrac{x}{1}=\dfrac{y-1}{1}=\dfrac{z-2}{1}$ và tiếp xúc với hai mặt  phẳng $(P)\colon 2x-z-4=0$, $(Q)\colon x-2y-2=0$. Tổng $P=a+b+c$ bằng bao nhiêu?
	 \shortans{6}
	 \loigiai{
		Gọi tâm mặt cầu là $I(a;a+1;a+2)\in  d$.\\
		Do mặt  cầu tiếp xúc với $(P)$ và $(Q)$ nên $\mathrm{d}\left(I,(P)\right)=\mathrm{d}\left(I,(Q)\right)$.\\
		Suy ra $\dfrac{\left| 2a-a-2-4 \right|}{\sqrt{5}}=\dfrac{\left| a-2(a+1)-2  \right|}{\sqrt{5}}\Leftrightarrow \left| a-6 \right|=\left| -a-4 \right|\Leftrightarrow a=1$.\\
		Tâm  mặt cầu là $I(1;2;3)$, suy ra $P=a+b+c=1+2+3=6$.
	} 
\end{bt}

% nguồn: Phan2_Lop_12_Hinh_hoc_Chuong_5
\begin{bt}%[2H5V3-2]%[ma:68fdf00f]
	Trong không gian $Oxyz$, cho hai đường thẳng $\Delta_1\colon \dfrac{x-4}{3}=\dfrac{y-1}{-1}=\dfrac{z+5}{-2}$ và $\Delta_2\colon \dfrac{x-2}{1}=\dfrac{y+3}{3}=\dfrac{z}{1}$. Trong tất cả mặt cầu tiếp xúc với cả hai đường thẳng $\Delta_1$ và $\Delta_2$. Gọi $(S)$ là mặt cầu có bán kính nhỏ nhất. Diện tích của mặt cầu đó có kết  quả bằng bao nhiêu (làm tròn đến hàng phần chục)?
	\shortans{18,8} 
	\loigiai{ 
		Ta có $\Delta_1$ đi qua $M(4;1;-5)$ và có vectơ chỉ  phương $\overrightarrow{u}_{{\Delta_1}}=(3;-1;-2)$.\\
		Ta có  $\Delta_2$ đi qua $N(2;-3;0)$ và có vectơ chỉ phương $\overrightarrow{u}_{{\Delta_2}}=(1;3;1)$.\\
		Ta có $\left[\overrightarrow{u}_{\Delta_1},\overrightarrow{u}_{\Delta_2} \right]=(5;-5;10)$.\\ 
		Vectơ $\overrightarrow{MN}=(-2;-4;5)$.\\ 
		Ta có $\left[\overrightarrow{u}_{\Delta_1},\overrightarrow{u}_{\Delta_2} \right]\cdot\overrightarrow{MN}=5\cdot(-2)-5\cdot(-4)+10\cdot5=60\ne 0$.\\ 
		Do đó $\Delta_1$ và $\Delta_2$ chéo nhau.\\ 
		Vì  mặt cầu $(S)$ tiếp xúc với hai đường thẳng $\Delta_1$, $\Delta_2$ và có bán kính nhỏ nhất, nên mặt cầu $(S)$ là mặt cầu đường kính $AB$ (với $AB$ là đoạn vuông góc chung của $\Delta_1$, $\Delta_2$).
		$$\mathrm{d}\left(\Delta_1,\Delta_2\right)=\dfrac{\left| \left[\overrightarrow{u}_{\Delta_1},\overrightarrow{u}_{\Delta_2} \right]\cdot\overrightarrow{MN} \right|}{\left| \left[\overrightarrow{u}_{\Delta_1},\overrightarrow{u}_{\Delta_2} \right] \right|}\dfrac{60}{\sqrt{5^2+(-5)^2+10^2}}=2\sqrt{6}.$$ 
		Vậy bán kính $R=\dfrac{AB}2=\dfrac{\mathrm{d}\left(\Delta_1,\Delta_2\right)}2=\sqrt{6}$.\\ 
		Diện tích mặt cầu  $(S)$ là $S=\pi R^2=6\pi \approx 18{,}8$.
	 }
\end{bt}

% nguồn: Phan2_Lop_12_Hinh_hoc_Chuong_5
\begin{bt}%[2H5V3-2]%[ma:4b81adaa]
	Trong không gian $Oxyz$, cho mặt cầu $(C)\colon (x-1)^2+(y-2)^2+(z+3)^2=11$. Đường thẳng $d$ là tiếp tuyến của mặt cầu $(C)$ tại điểm $A(2;1;0)$ có vectơ chỉ phương  là $\overrightarrow{u}_{d}=(-1;b;c)$ và vuông góc với đường thẳng $\Delta\colon\dfrac{x+1}{1}=\dfrac{y}2=\dfrac{z-2}{1}$. Tính giá trị $b+c$ (kết quả làm tròn đến hàng phần trăm).
	 \shortans{0,71}
	 \loigiai{
		 Mặt cầu $(C)$ có tâm $I=(1;2;-3)$, khi đó $\overrightarrow{IA}=(1;-1;3)$.\\
		Đường thẳng $\Delta$ có vectơ chỉ phương  $\overrightarrow{u}_{\Delta}=(1;2;1)$.\\
		Do $d$ là tiếp tuyến  của mặt cầu tại điểm $A(2;1;0)$ nên $\overrightarrow{IA}\cdot\overrightarrow{u}_d=0\Leftrightarrow -b+3c-1=0$.\\
		Lại có $d$ vuông góc với đường thẳng $\Delta $ nên  $\overrightarrow{u}_{\Delta }.\overrightarrow{u}_{d}=0\Leftrightarrow 2b+c-1=0$.\\
		Ta có hệ phương trình $\heva{& -b+3c-1=0 \\& 2b+c-1=0}\Leftrightarrow  \heva{
		& b=\dfrac2{7} \\&  c=\dfrac{3}{7}.}$\\
		Vậy  $b+c=\dfrac{5}{7}\approx 0{,}71$.
	} 
\end{bt}

% nguồn: Phan2_Lop_12_Hinh_hoc_Chuong_5
\begin{bt}%[2H5V3-2]%[ma:764c8127]
	Trong không gian với hệ trục tọa độ $Oxyz$, cho mặt phẳng $(P)\colon2x-y+2z-14=0$ và mặt cầu $(S)\colon x^2+y^2+z^2-2x+4y+2z-3=0$. Lấy $M\left(a;b;c\right)$ thuộc mặt cầu $(S)$ sao cho khoảng cách từ $M$ đến mặt phẳng $(P)$ là lớn nhất. Giá trị của  biểu thức $K=a+b+c$ là
	 \shortans{-5}
	\loigiai{ 
		Mặt cầu $(S)$ có tâm $I\left(1;-2;-1\right)$,  bán kính $R=3$.\\
		Ta  có $\mathrm{d}\left(I,(P)\right)$ $=\dfrac{\left| 2\cdot1-\left(-2\right)+2\cdot\left(-1\right)-14\right|}{\sqrt{2^2+\left(-1\right)^2+2^2}}=4>R$.\\
		Suy ra  mặt phẳng $(P)$ và mặt cầu $(S)$ không có điểm chung.\\
		Từ đó, điểm thuộc mặt cầu có khoảng cách nhỏ nhất hoặc lớn nhất  tới mặt phẳng $(P)$ là giao điểm của mặt cầu với đường thẳng qua $I$ và vuông góc với $(P)$.\\
		Trước hết ta lập phương trình đường thẳng $d$ đi qua $I$ và vuông  góc với $(P)$.\\
		Mặt phẳng  $(P)$ có véctơ pháp tuyến là $\overrightarrow{n}=\left(2;-1;2\right)$.\\
		Vì $d\perp(P)$ nên nhận $\overrightarrow{n}=\left(2;-1;2\right)$ làm véctơ chỉ phương.\\ 
		Từ đó  $d$ có phương trình $\heva{
		&  x=1+2t\\
		& y=-2-t\\ 
		&  z=-1+2t}$ với $\left(t\in\mathbb{R}\right)$.\\
		Ta tìm giao điểm của $d$ và $(S)$.\\ Xét  hệ phương trình $\heva{
		& x=1+2t\\ 
		& y=-2-t\\ 
		&  z=-1+2t\\
		 &{x^2}+y^2+z^2-2x+4y+2z-3=0}$\\
		$\Leftrightarrow\heva{ 
		&  x=1+2t\\
		 & y=-2-t\\
		 & z=-1+2t\\
		 &{\left(1+2t\right)^2}+\left(-2-t\right)^2+\left(-1+2t\right)^2-2\left(1+2t\right)+4\left(-2-t\right)+2\left(-1+2t\right)-3=0}$\\
		 $\Leftrightarrow\heva{&x=1+2t\\
		&  y=-2-t\\
		&  z=-1+2t\\
		&  9t^2-9=0}$ $\Leftrightarrow\hoac{
		&\heva{ 
		 & t=1\\
		&  x=3\\
		&  y=-3\\
		& z=1}\\ 
		&\heva{ 
		& t=-1\\ 
		&  x=-1\\
		 & y=-1\\
		&  z=-3.}}$\\
		Suy ra có hai giao điểm là  $A\left(3;-3;1\right)$ và $B\left(-1;-1;-3\right)$.\\
		Ta  có $\mathrm{d}\left(A,(P)\right)=\dfrac{\left| 2\cdot3-\left(-3\right)+2\cdot1-14\right|}{\sqrt{2^2+\left(-1\right)^2+2^2}}=1$; $\mathrm{d}\left(B,(P)\right)=\dfrac{\left| 2\cdot\left(-1\right)-\left(-1\right)+2\cdot\left(-3\right)-14\right|}{\sqrt{2^2+\left(-1\right)^2+2^2}}=7$.\\
		Suy ra $M\equiv B\left(-1;-1;-3\right)$. Từ đó $a=-1$;  $b=-1$; $c=-3$.\\
	Vậy  $K=-5$.}
\end{bt}

% nguồn: Phan2_Lop_12_Hinh_hoc_Chuong_5
\begin{bt}%[2H5V3-2]%[ma:33d31d8e]
	Trong không gian $Oxyz$, cho $A(-2;0;0)$; $B(0;-2;0)$; $C(0;0;-2)$. $D$ là điểm khác $O$ sao cho $DA$, $DB$, $DC$ đôi một vuông góc. $I(a;b;c)$ là tâm mặt cầu  ngọai tứ diện $ABCD$. Tính $S=a+b+c$.
	 \shortans{-1}
	 \loigiai{
		Gọi $D(x;y;z)$, suy ra $\overrightarrow{AD}=(x+2;y;z)$; $\overrightarrow{BD}=(x;y+2;z)$; $\overrightarrow{CD}=(x;y;z+2)$.\\ 
		Vì $DA$, $DB$, $DC$ đôi một vuông góc nên 
		\allowdisplaybreaks 
		\begin{eqnarray*}
			&&\heva{& \overrightarrow{AD}\cdot\overrightarrow{BD}=0 \\ & \overrightarrow{AD}\cdot\overrightarrow{CD}=0\\ & \overrightarrow{BD}\cdot\overrightarrow{CD}=0}\\&\Leftrightarrow&  \heva{
			& x(x+2)+y(y+2)+z^2=0 \\ 
			&  x(x+2)+y^2+z(z+2)=0 \\
			& x^2+y(y+2)+z(z+2)=0}\\ 
			&\Leftrightarrow&  \heva{
			 &x^2+y^2+z^2+2x+2y=0 \\
			& x^2+y^2+z^2+2x+2z=0 \\ 
			 &x^2+y^2+z^2+2y+2z=0}\\
			 &\Leftrightarrow& \heva{
			 &x=y=z\\
			 & x^2+x^2+x^2+2x+2x=0}\\
			&\Leftrightarrow&  \heva{
			 &x=y=z\\
			 & 3x^2+4x=0}\\
			&\Leftrightarrow& \heva{ 
			&x=y=z\\ 
			& \hoac{&x=0\text{ (loại)}\\&x=-\dfrac{4}{3}}}\\ 
			&\Leftrightarrow&  x=y=z=-\dfrac{4}{3}.
		\end{eqnarray*}
		$I(a;b;c)$ là tâm mặt cầu ngọai tiếp tứ diện  $ABCD$ nên
		\allowdisplaybreaks 
		\begin{eqnarray*}
			&&\heva{& IA=IB\\& IA=IC\\& IA=ID} 
			\Leftrightarrow \heva{ 
			 & (a+2)^2+b^2+c^2=a^2+(b+2)^2+c^2 \\
			&(a+2)^2+b^2+c^2=a^2+b^2+(c+2)^2  \\
			&(a+2)^2+{{b}^2}+c^2=\left(a+\dfrac{4}{3} \right)^2+\left(b+\dfrac{4}{3} \right)^2+\left(c+\dfrac{4}{3} \right)^2.}\\ 
			&\Leftrightarrow& \heva{& a=b \\& a=c \\& 4a+4=8a+\dfrac{16}{3}}\Leftrightarrow a=b=c=-\dfrac{1}{3}. 
		\end{eqnarray*}
		Vậy $a+b+c=-1$. 
	} 
\end{bt}

% nguồn: Phan2_Lop_12_Hinh_hoc_Chuong_5
\begin{bt}%[2H5V3-2]%[ma:f769d172]
	Trong không gian $Oxyz$, cho lăng trụ đứng $ABC.A'B'C'$ có $A(a; 0; 0),  B(-a; 0; 0),$\\
	$ C(0; 1; 0)$ và $B'(-a; 0; b)$, trong đó $a, b$ là các số thực dương và thỏa mãn $a+b=4$. Gọi $(\alpha)$ là mặt phẳng chứa $B'C$ và song song với $AC'$ sao khoảng cách từ $A$ đến $(\alpha)$ là lớn nhất. Khi đó, mặt cầu đi  qua sáu điểm $A$, $B$, $C$, $A'$, $B'$, $C'$ có bán kính $R$ bằng bao nhiêu? Làm tròn kết quả đến hàng phần trăm.
	\shortans{2,69} 
	\loigiai{ 
		\begin{center}
			\begin{tikzpicture}[scale=1,>=stealth,font=\footnotesize,line join=round,line cap=round]
					 \def\a{4}
					 \def\h{4.5}
					\path (0:0) coordinate (A) 
						++(0:1.3*\a) coordinate (C) 
						++(-150:3*\a/4) coordinate  (B)
						($(A)+(90:\h)$) coordinate  (A')
						($(B)+(90:\h)$)  coordinate (B')
						($(C)+(90:\h)$) coordinate (C') 
						(intersection of  A--C' and B'--C) coordinate (M)
						 ($(M)+(-90:\h/3)$) coordinate (I)
						($(I)+(-90:2)$) coordinate  (V);
					\draw[dashed]  (C')--(A)--(C)
						(A) arc (180:0:{2.6}  and {1.5})
						(M)--++(0,1.6) 
						(M)--(I)--(V)--(B)--(I); 
					\draw (C)--(C') 
						(B)--(B')--(C) 
						(A)--(A') 
						 (A)--(B)--(C)
						 (A')--(B')--(C')--cycle
						(A) arc  (-180:0:{2.6} and {1.5})
						 pic[draw,angle radius=2mm]{right angle=B--V--I};
					\foreach \x/\g  in {A/180,B/-45,C/0,A'/180,B'/90,C'/0,I/0}
					\fill (\x)  circle (1pt)
						 ($(\g:4mm)+(\x)$) node {$\x$};
				\end{tikzpicture}
		\end{center}
		Ta có $A(a; 0; 0), C'(0; 1; b)  \Rightarrow \overrightarrow{AC'}=(-a; 1; b)$;\\
		$B'(-a; 0; b), C(0; 1; 0) \Rightarrow \overrightarrow{B'C}=(a; 1;  -b)$.\\
		Vì $(\alpha)$ chứa $B'C$ và song song với $A'C$ nên $(\alpha)$ có một vectơ  pháp tuyến là $\overrightarrow{n} =\left[\overrightarrow{AC'},\overrightarrow{B'C}\right]=(-2b; 0; -2a)$.\\
		Phương trình mặt phẳng $(\alpha)$ là $bx + az =0$  với $a+b=4$.\\
		Khoảng cách từ điểm $A$ đến mặt phẳng $(\alpha)$  là $d\left(A, (\alpha)\right)=\dfrac{ab}{\sqrt{a^2+b^2}}=\dfrac{-a^2+4a}{\sqrt{2a^2-8a+16}}$.\\
		Trên $[0;4]$, hàm số $f(x)=\dfrac{-x^2+4x}{\sqrt{2x^2-8x+16}}$  đạt giá trị lớn nhất bằng $\sqrt{2}$ tại $x=2$.\\
		Suy  ra $a=b=2$.\\
		Khi đó $A(2; 0; 0), B(-2; 0; 0), C(0; 1;  0), B'(-2; 0; 2)$.\\
		Do đó $AC=BC=\sqrt{5}$ nên tam giác $ABC$ cân tại $C$ và $OC  \perp AB$.\\
		Phương trình tham số của đường thẳng $CO \colon \heva{&x=0\\&y=t\\&z=0}  \left(t \in \mathbb{R}\right)$.\\
		Diện tích tam giác $ABC$  là $S=\dfrac{1}{2}\cdot AB \cdot CO= \dfrac{1}{2}\cdot4 \cdot 1=2$.\\
		Gọi $r$ là bán kính đường  tròn ngoại tiếp tam giác $ABC$, ta có
		\[r=\dfrac{AB \cdot BC  \cdot CA}{4S}=\dfrac{4\cdot \sqrt{5} \cdot \sqrt{5}}{4 \cdot 2}=\dfrac{5}{2}.\]
		Gọi $I$ là tâm đường tròn ngoại tiếp  tam giác $ABC$, ta có $IA^2=IC^2 \Rightarrow t=-\dfrac{3}{2}$.\\
		 Do đó $I\left(0; -\dfrac{3}{2}; 0\right)$.\\
		Lăng trụ đứng $ABC.A'B'C'$  có chiều cao $BB'=2$.\\
		Gọi $I'$  là tâm đường tròn ngoại tiếp tam giác $A'B'C'$ thì $I'\left(0; -\dfrac{3}{2}; 2\right)$.\\
		Khi đó trung điểm $J\left(0;  -\dfrac{3}{2}; 1\right)$ của $II'$ là tâm mặt cầu ngoại tiếp lăng trụ $ABC.A'B'C'$.\\
		Bán kính mặt cầu đi qua sáu điểm $A$, $B$, $C$, $A'$, $B'$, $C'$  là $R=JC=\dfrac{\sqrt{29}}{2} \approx 2{,}69$.
	} 
\end{bt}

% nguồn: Phan2_Lop_12_Hinh_hoc_Chuong_5
\begin{bt}%[2H5V3-2]%[ma:8459d640]
	Trong không gian với hệ tọa độ $Oxyz$, cho đường thẳng $d\colon \dfrac{x-1}{2}=\dfrac{y-1}{-1}=\dfrac{z-4}{2}$ và mặt cầu $(S)\colon x^2+y^2+z^2=9$. Gọi $(P)$ là mặt phẳng  có phương trình dạng: $ax+by+cz+9=0$. Biết rằng $(P)$ chứa $d$ và $(P)$ tiếp xúc với $(S)$. Tính giá trị biểu thức $a+b+c$.
	\shortans{-3} 
	\loigiai{ 
		Mặt cầu $(S)$ có tâm $I(0,0,0)$ và bán kính $R=3$.\\ 
		 Đường thẳng $d$ đi qua điểm $M(1,1,4)$ và có véc-tơ chỉ phương $\overrightarrow{u_d}=(2,-1, 2)$.\\
		Mặt phẳng  $(P)\colon ax+by+cz+9=0$ có véc-tơ pháp tuyến $\overrightarrow{n_P}=(a,b,c)$.\\
		Vì $(P)$ chứa $d$ nên $M  \in(P)$ và $\overrightarrow{n_P}\perp \overrightarrow{u_d}$.\\
		Ta có $M \in(P)\Rightarrow a\cdot 1+b\cdot 1+c\cdot 4+9=0 \Rightarrow  a+b+4c+9=0$.\hfill(1)\\
		Lại có $\overrightarrow{n_P}\cdot \overrightarrow{u_d}=0 \Rightarrow 2a-b+2c=0  \Rightarrow b=2a+2c$.\hfill (2)\\
		Thay (2) vào (1):  $a+(2a+2c)+4c+9=0 \Rightarrow a=-2c-3$.\hfill (3)\\
		Từ (2) và (3), suy ra $b=2(-2c-3)+2c=-4c-6+2c=-2c-6$.\\ 
		 Vậy $\overrightarrow{n_P}=(-2c-3,-2c-6, c)$.\\
		 Mặt khác $(P)$ tiếp xúc với $(S)$ nên $\mathrm{d}(I,(P))=R$.\\
		Do đó 
		 \allowdisplaybreaks
		\begin{eqnarray*}
			&&\dfrac{|a\cdot 0+b\cdot0 +c \cdot 0+9|}{\sqrt{a^2+b^2+c^2}}=3 \\ 
			&\Leftrightarrow&  a^2+b^2+c^2=9.\\
			&\Leftrightarrow&(-2c-3)^2+(-2c-6)^2+c^2=9.\\ 
			&\Leftrightarrow&(4c^2+12c+9)+(4c^2+24c+36)+c^2=9.\\ 
			&\Leftrightarrow&9c^2+36c+45=9.\\ 
			 &\Leftrightarrow&9c^2+36c+36=0.\\
			&\Leftrightarrow&c^2+4c+4=0.\\ 
			&\Leftrightarrow&  c=-2.
		\end{eqnarray*}
		Khi đó  $a=-2\cdot (-2)-3=4-3=1$.\\
		$b=-2\cdot (-2)-6=4-6=-2$.\\ 
		 Vậy $a=1, b=-2, c=-2$.\\
		$a+b+c=1+(-2)+(-2)=-3$. 
	} 
\end{bt}

% nguồn: Phan2_Lop_12_Hinh_hoc_Chuong_5
\begin{bt}%[2H5C3-2]%[ma:b874802b]
	Trong không gian $Oxyz$, cho đường thẳng $\Delta \colon\dfrac{x-2}{-3}=\dfrac{y-1}{2}=\dfrac{z-1}{2}$, $(P)\colon x+2y-2z-2=0$, $(Q)\colon x+2y-2z+4=0$. Gọi mặt cầu $S\left(I,R\right)$ có tâm $I$ thuộc $\Delta$ và tiếp xúc với $(P)$,  $(Q)$. Khi đó đường kính của mặt cầu có giá trị bằng bao nhiêu?
	\shortans{2} 
	\loigiai{ 
		Ta có phương trình tham  số của $\Delta $ là $\heva{
		&  x=2-3t\\
		&  y=1+2t\\
		& z=1+2t.}$\\ 
		Vì  $I\in\Delta\Rightarrow I\left(2-3t;1+2t;1+2t\right)$.\\
		Mặt cầu $(S)$ có tâm $I\left(2-3t;1+2t;1+2t\right)$ thuộc $\Delta $ và tiếp xúc với $(P)$, $(Q)$ nên ta  có\\
		$\mathrm{d}\left(I,(P)\right)=\mathrm{d}\left(I,(Q)\right)$  $\Leftrightarrow\dfrac{\left|\left(2-3t\right)+2\left(1+2t\right)-2\left(1+2t\right)-2\right|}{\sqrt{1^2+2^2+\left(-2\right)^2}}=\dfrac{\left|\left(2-3t\right)+2\left(1+2t\right)-2\left(1+2t\right)+4\right|}{\sqrt{1^2+2^2+\left(-2\right)^2}}$\\
		$\Leftrightarrow\left|-3t\right|=\left|-3t+6\right|\Rightarrow  t=1$.\\
		Suy ra mặt cầu $(S)$  có tâm $I\left(-1;3;3\right)$ và bán kính $R=1$.\\
	Vậy đường kính của mặt cầu $(S)$ bằng $2$.} 
\end{bt}


%%%%%%%%%%%%%%%% TỰ LUẬN %%%%%%%%%%%%%%%%

% nguồn: Ok_TuanLuan_Toan11_Lan2
\begin{ex}%[2H5V3-2]%[ma:dbb27e8a]
	Trong không gian $Oxyz$, cho mặt cầu $(S)\colon (x-1)^2+(y-2)^2+(z+3)^2=4$ và mặt phẳng $(P)\colon x+y-z+4=0$. Xét  vị trí tương đối của mặt phẳng $(P)$ và mặt cầu $(S)$.
	 \loigiai{
		Mặt cầu $(S)$ có tâm $I(1;2;-3)$, bán kính $R=\sqrt{4}=2$  và $\mathrm{d}\left(I,(P)\right)=\dfrac{\left| 1+2+3+4 \right|}{\sqrt{1^2+1^2+(-1)^2}}=\dfrac{10}{\sqrt{3}}>R$.\\
		Vậy mặt phẳng $(P)$ và mặt  cầu $(S)$ không có điểm chung.
	 }
\end{ex}

% nguồn: Ok_TuanLuan_Toan11_Lan2
\begin{ex}%[2H5V3-2]%[ma:1e03a3c0]
	Trong không gian $Oxyz$, đường thẳng $d\colon\heva{& x=4+2t \\& y=3+t \\& z=3+2t}$, $t\in \mathbb{R}$. Viết phương trình mặt cầu $(S)$ có tâm $I(-2;3;0)$ và đường thẳng  $d$ cắt mặt cầu $(S)$ tại hai điểm phân biệt $A$, $B$ sao cho độ dài đọan thẳng $AB=8$.
	 \loigiai{
		\begin{center}
			\begin{tikzpicture}[scale=1,font=\footnotesize,line join=round,line cap=round,>=stealth]
					\def \r{2} 
					\path 
						(0,0)  coordinate (I)
						 (160:\r) coordinate (A)
						(60:\r) coordinate (B) 
						 ($(A)!(I)!(B)$) coordinate (H)
						(180:\r) coordinate  (N)
						(intersection of  A--B and I--N) coordinate (M) node[above]{$d$};
					\draw (I) circle  (\r);
					\draw (I)--(H) 
						 (M)--(B)
						 (A)--(I)--(B);
					 \pic[draw,angle radius=2mm] {right angle = I--H--B};
					\foreach \x/\g in {A/150,B/60,I/-90,H/90}\fill (\x)  circle (1pt)+(\g:.3) node{$\x$};
				\end{tikzpicture}
		\end{center}
		Gọi $H$ là trung điểm của  $AB$, suy ra $IH\perp AB$, $AH=4$.\\
		Đường thẳng $d$ đi qua điểm $M(4;3;3)$  và có $1$ vectơ chỉ phương $\overrightarrow{u}=(2;1;2)$.\\
		Ta có  $\overrightarrow{IM}=(6;0;3)\Rightarrow \left[\overrightarrow{IM},\overrightarrow{u} \right]=(-3;-6;6)$.\\
		$\Rightarrow IH=\mathrm{d}\left(I,d \right)=\dfrac{\left| \left[\overrightarrow{IM},\overrightarrow{u} \right] \right|}{\left| \overrightarrow{u} \right|}=\dfrac{\sqrt{(-3)^2+(-6)^2+6^2}}{\sqrt{2^2+1^2+2^2}}=3$.\\ 
		Khi  đó $R^2=IA^2=IH^2+AH^2=25\Rightarrow R=5$.\\
		Mặt cầu $(S)$ có tâm $I(-2;3;0)$ và  bán kính $R=5$ nên có phương trình là
		 $(x+2)^2+(y-3)^2+z^2=25$.
	 }
\end{ex}

% nguồn: Ok_TuanLuan_Toan11_Lan2
\begin{ex}%[2H5V3-2]%[ma:1f33257b]
	Trong không gian $Oxyz$, cho mặt phẳng $(P)\colon 2x+3y+z+17=0$ và mặt cầu $(S)\colon x^2+y^2+z^2-2x+4y-2z-8=0$. Xét vị trí tương đối của mặt  phẳng $(P)$ và mặt cầu $(S)$.
	\loigiai{ 
		Mặt  cầu $(S)$ có tâm $I(1;-2;1)$, bán kính $R=\sqrt{14}$ và $\mathrm{d}\left(I,(P) \right)=\dfrac{\left| 2-6+1+17 \right|}{\sqrt{2^2+3^2+1^2}}=\sqrt{14}$.\\
		Vậy mặt phẳng $(P)$ tiếp xúc với  mặt cầu $(S)$.
	} 
\end{ex}

% nguồn: Ok_TuanLuan_Toan11_Lan2
\begin{ex}%[2H5V3-2]%[ma:b4e5fa99]
	Trong không gian $Oxyz$, cho mặt cầu $(S)\colon (x+1)^2+(y-2)^2+z^2=9$ và đường thẳng có  phương trình $\Delta\colon \dfrac{x-1}2=\dfrac{y-2}{-2}=\dfrac{z-3}{1}$. Tìm vị trí tương đối của đường thẳng $\Delta$ và mặt cầu $(S)$.
	\loigiai{ 
		Mặt cầu $(S)\colon (x+1)^2+(y-2)^2+z^2=9$ có tâm $I(-1;2;0)$ bán kính $R=3$.\\ 
		Gọi $H$ là hình  chiếu vuông góc của $I$ trên $\Delta$ thì ta có $H(1+2t;2-2t;3+t)$, $t\in \mathbb{R}$.\\
		Khi đó $\overrightarrow{IH}=(2+2t;-2t;3+t)$. Một vectơ chỉ  phương của $\Delta$ là $\overrightarrow{u}=(2;-2;1)$.\\
		Ta có $\overrightarrow{IH}\cdot\overrightarrow{u}=0\Leftrightarrow 4+4t+4t+3+t=0\Leftrightarrow  t=-\dfrac{7}{9}$.\\
		$\overrightarrow{IH}=\left( \dfrac{4}{9};\dfrac{14}{9};\dfrac{20}{9} \right)\Rightarrow IH=\sqrt{\dfrac{16}{81}+\dfrac{196}{81}+\dfrac{400}{81}}=\dfrac{2\sqrt{17}}{3}$.\\ 
		Suy  ra $\mathrm{d}\left(I,\Delta\right)=\dfrac{2\sqrt{17}}{3}<R$.\\
		Vậy đường thẳng $\Delta$ cắt mặt cầu  $(S)$ tại hai điểm phân biệt.
	 }
\end{ex}

% nguồn: Ok_TuanLuan_Toan11_Lan2
\begin{ex}%[2H5V3-2]%[ma:2431641c]
	Trong không gian $Oxyz$, cho hai đường thẳng chéo nhau $d\colon \dfrac{x-2}2=\dfrac{y-3}{3}=\dfrac{z+4}{-5}$ và $d'\colon \dfrac{x+1}{3}=\dfrac{y-4}{-2}=\dfrac{z-4}{-1}$. Viết phương trình mặt cầu có bán kính nhỏ nhất tiếp xúc với cả hai đường  thẳng $d$, $d'$.
	 \loigiai{
		Đường  thẳng $d$ có một vectơ chỉ phương là $\overrightarrow{u}_{d}=(2;3;-5)$.\\
		Đường thẳng $d'$ có một vectơ  chỉ phương là $\overrightarrow{u}_{d'}=(3;-2;-1)$.\\
		Gọi $M\in d\Rightarrow M(2+2m;3+3m;-4-5m)$ và $N\in  {d}'\Rightarrow N(-1+3n;4-2n;4-n)$.\\
		Khi đó ta  có $\overrightarrow{MN}=(-3+3n-2m;1-2n-3m;8-n+5m)$.\\
		$MN$ là đường vuông góc chung của $d$ và ${d}'$ khi  và chỉ khi
		 \allowdisplaybreaks
		\begin{eqnarray*}
			&&\heva{& MN\perp d \\& MN\perp d'}\Leftrightarrow \heva{&\overrightarrow{MN}\cdot\overrightarrow{u}_{d}=0 \\  & \overrightarrow{MN}\cdot\overrightarrow{u}_{d'}=0}\\
			&\Leftrightarrow& \heva{& 2(-3+3n-2m)+3(1-2n-3m)-5(8-n+5m)=0\\&  3(-3+3n-2m)-2(1-2n-3m)-1(8-n+5m)=0}\\
			&\Leftrightarrow& \heva{& -38m+5n=43  \\
			& -5m+14n=19}\Leftrightarrow \heva{& m=-1  \\& n=1.}
		\end{eqnarray*}
		Suy ra  $M(0;0;1)$, $N(2;2;3)$.\\
		Gọi $I$ là  trung điểm của $MN$, suy ra $I(1;1;2)$ và $IM=IN=\sqrt{1^2+1^2+1^2}=\sqrt{3}$.\\
		Khi đó mặt cầu tâm $I$, bán kính $R=\sqrt{3}$ tiếp xúc  với hai đường thẳng $d$, $d'$ lần lượt tại $M$ và $N$.\\
		Do $MN$ là khoảng cách giữa hai đường thẳng $d$, $d'$ nên mặt cầu đường kính $MN$ là mặt cầu có bán kính nhỏ nhất tiếp xúc với hai đường thẳng $d$, $d'$.\\ 
		Vậy phương trình mặt  cầu cần tìm là $(x-1)^2+(y-1)^2+(z-2)^2=3$.
	} 
\end{ex}

% nguồn: Ok_TuanLuan_Toan11_Lan2
\begin{ex}%[2H5V3-2]%[ma:e2bdba25]
	Trong không gian với hệ tọa độ $Oxyz$ cho mặt cầu $(S)\colon x^2+y^2+z^2-4x+2y-2z-3=0$ và điểm $A(5;3;-2)$. Một đường thẳng $d$ thay đổi luôn đi qua $A$ và luôn cắt mặt cầu tại hai điểm phân biệt $M$, $N$. Tính  giá trị nhỏ nhất của biểu thức $S=AM+4AN$.
	\loigiai{ 
		\begin{center}
			\begin{tikzpicture}[scale=1,font=\footnotesize,line join=round,line cap=round,>=stealth]
					\def  \r{2}
					 \path
						(0,0)  coordinate (I)
						(160:\r)  coordinate (N)
						(60:\r)  coordinate (M)
						 ($(M)!(I)!(N)$) coordinate (H)
						(180:\r) coordinate  (B)
						(intersection of M--N  and I--B) coordinate (A);
					\draw  (I) circle (\r);
					\draw  (I)--(A)--(M)
						(I)--(H); 
					\pic[draw,angle radius=2mm] {right  angle = I--H--N};
					\foreach \x/\g in {A/100,N/120,M/50,I/-90,H/90}\fill (\x) circle (1pt)+(\g:.3)  node{$\x$};
				\end{tikzpicture}
		\end{center}
		Mặt cầu $(S)$ có tâm  $I(2;-1;1)$, bán kính $R=3$.\\
		$AI=\sqrt{34}>R$, suy ra $A$ nằm ngoài mặt cầu $(S)$.\\ 
		Do hai điểm $M$, $N$ nằm ở vị trí hai đầu một dây cung nên  để $S_{\min}$ thì $N$ nằm giữa $A$ và $M$.\\
		Gọi  $H$ là trung điểm $MN$, suy ra $IH\perp MN$, $NH=\dfrac{1}2MN$.\\
		 $S=AH+NH+4(AH-NH)=5AH-3NH$\\
		 $S=5\sqrt{AI^2-IH^2}-3\sqrt{R^2-IH^2}=5\sqrt{34-x^2}-3\sqrt{9-x^2}$, với $x=IH$.\\
		 Xét hàm số $f(x)=5\sqrt{34-x^2}-3\sqrt{9-x^2}$ ($0\le x<3$).\\
		Đạo hàm $f'(x)=\dfrac{-5x}{\sqrt{34-x^2}}+\dfrac{3x}{\sqrt{9-x^2}}=x\left( \dfrac{-5}{\sqrt{34-x^2}}+\dfrac{3}{\sqrt{9-x^2}}  \right)$.\\
		Xét 
		\allowdisplaybreaks 
		\begin{eqnarray*}
			&&  \dfrac{-5}{\sqrt{34-x^2}}+\dfrac{3}{\sqrt{9-x^2}}>0\\
			&\Leftrightarrow& 5\sqrt{9-x^2}<3\sqrt{34-x^2}\\ 
			&\Leftrightarrow& 225-25x^2<9\cdot 34-9x^2\Leftrightarrow  16x^2+81>0\text{ (luôn đúng)}.
		\end{eqnarray*}
		Suy  ra $f'(x)\ge 0,\forall x\in [0;3)$, $f'(x)=0\Leftrightarrow x=0\Rightarrow f(x)$ đồng biến trên $[0;3)$.\\
		Suy ra $\underset{[0;3)}{\mathop{\min  }}\,f(x)=f(0)=5\sqrt{34}-9$.
	 }
\end{ex}

% nguồn: Phan2_Lop_12_Hinh_hoc_Chuong_5
\begin{bt}%[2H5V3-2]%[ma:2434defc]
	Trong không gian $Oxyz$, cho mặt phẳng $(P)\colon 2x+3y+z+17=0$ và mặt cầu $(S)\colon x^2+y^2+z^2-2x+4y-2z-8=0$. Xét vị trí tương đối của mặt phẳng $(P)$ và mặt cầu  $(S)$.
	 \loigiai{
		Mặt cầu $(S)$ có tâm $I(1;-2;1)$, bán kính $R=\sqrt{14}$ và  $\mathrm{d}\left(I,(P) \right)=\dfrac{\left| 2-6+1+17 \right|}{\sqrt{2^2+3^2+1^2}}=\sqrt{14}$.\\
		Vậy mặt  phẳng $(P)$ tiếp xúc với mặt cầu $(S)$.
	 }
\end{bt}

% nguồn: Phan2_Lop_12_Hinh_hoc_Chuong_5
\begin{bt}%[2H5V3-2]%[ma:b17abcdd]
	Trong không  gian $Oxyz$, cho mặt cầu $(S)\colon (x-1)^2+(y-2)^2+(z+3)^2=4$ và mặt phẳng $(P)\colon x+y-z+4=0$. Xét vị trí tương đối của mặt phẳng $(P)$ và mặt cầu $(S)$.
	 \loigiai{
		Mặt cầu $(S)$ có tâm $I(1;2;-3)$, bán kính $R=\sqrt{4}=2$ và $\mathrm{d}\left(I,(P)\right)=\dfrac{\left| 1+2+3+4 \right|}{\sqrt{1^2+1^2+(-1)^2}}=\dfrac{10}{\sqrt{3}}>R$.\\ 
		Vậy mặt phẳng $(P)$ và mặt cầu $(S)$  không có điểm chung.
	} 
\end{bt}

% nguồn: Phan2_Lop_12_Hinh_hoc_Chuong_5
\begin{bt}%[2H5V3-2]%[ma:c7c7f192]
	Cho ba điểm  $A(1; 0; 0)$, $B(0; 2; 0)$ và $C(0; 0; 3)$. Chứng minh rằng nếu điểm $M(x; y; z)$ thoả mãn $M A^2=M B^2+M C^2$ thì $M$ thuộc một mặt cầu $(S)$. Tìm tâm và bán kính của $(S)$.
	\loigiai{ 
		 \allowdisplaybreaks
		\begin{eqnarray*}
			 \text{ Ta có}
			&&  M A^2=M B^2+M C^2\\
			 &\Leftrightarrow &\left(1-x \right)^2+y^2+z^2=x^2+\left(2-y \right)^2+z^2+x^2+y^2+\left(3-z \right)^2 .\\
			&\Leftrightarrow &  1-2x+x^2+y^2+z^2=x^2+4-4y+y^2+z^2+x^2+y^2+9-6z+z^2.\\
			&\Leftrightarrow & x^2+y^2+z^2+2x-4y-6z+12=0~ \text{ thoả  điều kiện}~ a^2+b^2+c^2-d=1+4+9-12=2>0.
		\end{eqnarray*}
		Vậy $M$ thuộc một mặt cầu $(S)$ có tâm $I\left( -1;2;3\right) $ và bán  kinh $R=\sqrt{2}$.
	} 
\end{bt}

% nguồn: Phan2_Lop_12_Hinh_hoc_Chuong_5
\begin{bt}%[2H5V3-2]%[ma:6a3ee938]
	Trong không gian $Oxyz$, cho mặt cầu $(S)\colon (x+1)^2+(y-2)^2+z^2=9$ và đường thẳng có phương trình $\Delta\colon \dfrac{x-1}2=\dfrac{y-2}{-2}=\dfrac{z-3}{1}$. Tìm vị trí tương đối của đường thẳng $\Delta$ và mặt  cầu $(S)$.
	\loigiai{ 
		 Mặt cầu $(S)\colon (x+1)^2+(y-2)^2+z^2=9$ có tâm $I(-1;2;0)$ bán kính $R=3$.\\
		Gọi $H$ là hình chiếu vuông góc của $I$ trên $\Delta$ thì ta có $H(1+2t;2-2t;3+t)$, $t\in \mathbb{R}$.\\ 
		Khi đó $\overrightarrow{IH}=(2+2t;-2t;3+t)$.  Một vectơ chỉ phương của $\Delta$ là $\overrightarrow{u}=(2;-2;1)$.\\
		Ta  có $\overrightarrow{IH}\cdot\overrightarrow{u}=0\Leftrightarrow 4+4t+4t+3+t=0\Leftrightarrow t=-\dfrac{7}{9}$.\\
		$\overrightarrow{IH}=\left( \dfrac{4}{9};\dfrac{14}{9};\dfrac{20}{9} \right)\Rightarrow IH=\sqrt{\dfrac{16}{81}+\dfrac{196}{81}+\dfrac{400}{81}}=\dfrac{2\sqrt{17}}{3}$.\\ 
		Suy  ra $\mathrm{d}\left(I,\Delta\right)=\dfrac{2\sqrt{17}}{3}<R$.\\
		Vậy đường thẳng $\Delta$ cắt mặt cầu  $(S)$ tại hai điểm phân biệt.
	} 
\end{bt}

% nguồn: Phan2_Lop_12_Hinh_hoc_Chuong_5
\begin{bt}%[2H5V3-2]%[ma:de0525d7]
	Trong không gian $Oxyz$, cho hai đường thẳng chéo nhau $d\colon \dfrac{x-2}2=\dfrac{y-3}{3}=\dfrac{z+4}{-5}$ và $d'\colon \dfrac{x+1}{3}=\dfrac{y-4}{-2}=\dfrac{z-4}{-1}$. Viết phương  trình mặt cầu có bán kính nhỏ nhất tiếp xúc với cả hai đường thẳng $d$, $d'$.
	 \loigiai{
		Đường thẳng  $d$ có một vectơ chỉ phương là $\overrightarrow{u}_{d}=(2;3;-5)$.\\
		Đường thẳng $d'$ có một  vectơ chỉ phương là $\overrightarrow{u}_{d'}=(3;-2;-1)$.\\
		Gọi $M\in d\Rightarrow M(2+2m;3+3m;-4-5m)$  và $N\in {d}'\Rightarrow N(-1+3n;4-2n;4-n)$.\\
		 Khi đó ta có $\overrightarrow{MN}=(-3+3n-2m;1-2n-3m;8-n+5m)$.\\
		$MN$ là đường  vuông góc chung của $d$ và ${d}'$ khi và chỉ khi
		 \allowdisplaybreaks
		\begin{eqnarray*}
			&&\heva{& MN\perp d \\& MN\perp d'}\Leftrightarrow \heva{&\overrightarrow{MN}\cdot\overrightarrow{u}_{d}=0 \\ &  \overrightarrow{MN}\cdot\overrightarrow{u}_{d'}=0}\\
			&\Leftrightarrow&  \heva{& 2(-3+3n-2m)+3(1-2n-3m)-5(8-n+5m)=0\\& 3(-3+3n-2m)-2(1-2n-3m)-1(8-n+5m)=0}\\
			&\Leftrightarrow& \heva{& -38m+5n=43  \\
			& -5m+14n=19}\Leftrightarrow \heva{& m=-1 \\& n=1.} 
		\end{eqnarray*}
		Suy ra  $M(0;0;1)$, $N(2;2;3)$.\\
		Gọi $I$ là trung điểm của $MN$, suy ra  $I(1;1;2)$ và $IM=IN=\sqrt{1^2+1^2+1^2}=\sqrt{3}$.\\
		Khi  đó mặt cầu tâm $I$, bán kính $R=\sqrt{3}$ tiếp xúc với hai đường thẳng $d$, $d'$ lần lượt tại $M$ và $N$.\\
		Do $MN$ là khoảng cách giữa hai đường thẳng $d$, $d'$ nên mặt cầu đường  kính $MN$ là mặt cầu có bán kính nhỏ nhất tiếp xúc với hai đường thẳng $d$, $d'$.\\
		Vậy phương  trình mặt cầu cần tìm là $(x-1)^2+(y-1)^2+(z-2)^2=3$.
	 }
\end{bt}

% nguồn: Phan2_Lop_12_Hinh_hoc_Chuong_5
\begin{bt}%[2H5V3-2]%[ma:fed92675]
	Trong không gian $Oxyz$, đường thẳng $d\colon\heva{& x=4+2t \\& y=3+t \\& z=3+2t}$, $t\in \mathbb{R}$. Viết phương trình mặt cầu $(S)$ có tâm $I(-2;3;0)$ và đường thẳng $d$ cắt mặt cầu $(S)$ tại hai điểm phân biệt $A$, $B$ sao cho độ  dài đọan thẳng $AB=8$.
	\loigiai{ 
		\begin{center}
			\begin{tikzpicture}[scale=1,font=\footnotesize,line join=round,line cap=round,>=stealth]
					 \def \r{2}
					 \path
						(0,0)  coordinate (I)
						(160:\r) coordinate (A) 
						(60:\r) coordinate (B) 
						 ($(A)!(I)!(B)$) coordinate (H)
						(180:\r)  coordinate (N)
						(intersection of A--B  and I--N) coordinate (M) node[above]{$d$};
					\draw  (I) circle (\r);
					\draw (I)--(H) 
						(M)--(B) 
						(A)--(I)--(B); 
					\pic[draw,angle radius=2mm]  {right angle = I--H--B};
					\foreach \x/\g in {A/150,B/60,I/-90,H/90}\fill (\x) circle  (1pt)+(\g:.3) node{$\x$};
				\end{tikzpicture}
		\end{center}
		Gọi $H$ là trung điểm của $AB$, suy ra $IH\perp  AB$, $AH=4$.\\
		Đường thẳng $d$ đi qua điểm $M(4;3;3)$ và có $1$  vectơ chỉ phương $\overrightarrow{u}=(2;1;2)$.\\
		Ta có  $\overrightarrow{IM}=(6;0;3)\Rightarrow \left[\overrightarrow{IM},\overrightarrow{u} \right]=(-3;-6;6)$.\\
		$\Rightarrow IH=\mathrm{d}\left(I,d \right)=\dfrac{\left| \left[\overrightarrow{IM},\overrightarrow{u} \right] \right|}{\left|  \overrightarrow{u} \right|}=\dfrac{\sqrt{(-3)^2+(-6)^2+6^2}}{\sqrt{2^2+1^2+2^2}}=3$.\\
		Khi đó $R^2=IA^2=IH^2+AH^2=25\Rightarrow R=5$.\\ 
		 Mặt cầu $(S)$ có tâm $I(-2;3;0)$ và bán kính $R=5$ nên có phương trình là
		$(x+2)^2+(y-3)^2+z^2=25$. 
	} 
\end{bt}

% nguồn: Phan2_Lop_12_Hinh_hoc_Chuong_5
\begin{bt}%[2H5V3-2]%[ma:c05e1ab7]
	Trong không gian với hệ tọa  độ $Oxyz$, cho hai điểm $A(3;-2;6)$, $B(0;1;0)$ và mặt cầu $(S)\colon (x-1)^2+(y-2)^2+(z-3)^2=25$. Viết phương trình mặt phẳng $(P)$ đi qua $A$, $B$ và cắt $(S)$ theo giao tuyến là đường tròn có bán kính nhỏ nhất.
	 \loigiai{
		\begin{center}
			\begin{tikzpicture}[scale=0.8,font=\footnotesize,line join=round,line cap=round,>=stealth]
					\def\R{3} 
					\def\g{30} 
					\pgfmathsetmacro\r{\R*cos  (\g)}
					\path 
						 (0,0) coordinate (I)
						(-\g:\R) coordinate (P) 
						 (-180+\g:\R) coordinate (M)
						($(M)!0.5!(P)$)  coordinate (H);
					 \draw
						(I)  circle (\R)
						(P) arc (0:-180: {\r} and {\r/4})  coordinate[pos=0.45](N);
					 \path
						 (N)++(-150:0.25*\R) coordinate (A)
						($(A)!3!(N)$) coordinate  (B)
						 ($(A)!0.6!(B)$) coordinate (K);
					\draw[dashed] 
						(P) arc  (0:180: {\r} and {\r/4})
						(I)--(M) node[midway,sloped,above]{$R$}--(H)  node[midway,sloped,below]{$r$}--(I)--(K)--(H)
						 (N)--($(A)!1.5!(B)$);
					\draw (N)--($(B)!1.5!(A)$); 
					 \pic[draw,angle radius=0.2cm]{right angle=I--H--M};
					\pic[draw,angle  radius=0.2cm]{right angle=B--K--I};
					\foreach \diem/\goc in {I/90,A/-90,B/-60,H/-90,K/0} \fill  (\diem) circle (1pt)
						($(\diem)+(\goc:3mm)$)  node{$\diem$};
				\end{tikzpicture}
		\end{center}
		Mặt cầu $(S)$ có tâm $I(1;2;3)$,  bán kính $R=5$.\\
		Gọi $H$,  $K$ lần lượt là hình chiếu của $I$ trên $(P)$ và $AB$.
		Bán kính  đường tròn giao tuyến là $r=\sqrt{R^2-IH^2}=\sqrt{25-IH^2}$
		$\Rightarrow r_{\min}\Leftrightarrow IH_{\max}$.\\ 
		Mà $IH\leq IK$ nên $IH_{\max}=IK\Leftrightarrow H=K\Leftrightarrow IK\perp(P)$.\\ 
		Ta có $\overrightarrow{AB}=(-3;3;-6)=-3(1;-1;2)\Rightarrow\overrightarrow{u}=(1;-1;2)$ là VTCP của đường  thẳng $AB$.\\
		$\Rightarrow$  phương trình tham số của $AB \heva{&x=t\\ &y=1-t\\ &z=2t}\Rightarrow K(t;1-t;2t)$.\\
		 $\Rightarrow\overrightarrow{IK}=(t-1;-t-1;2t-3)$.\\
		Vì $IK\perp AB$ nên $\overrightarrow{IK}\cdot \overrightarrow{u}=0\Leftrightarrow t-1+t+1+4t-6=0\Leftrightarrow t=1\Rightarrow \heva{& K(1;0;2)\\&\overrightarrow{IK}=(0;-2;-1).}$\\ 
		Suy ra  phương trình $(P)$ là $-2y-1(z-2)=0\Leftrightarrow 2y+z-2=0$.
	} 
\end{bt}

% nguồn: Phan2_Lop_12_Hinh_hoc_Chuong_5
\begin{bt}%[2H5V3-2]%[ma:cf6fbea5]
	Trong  không gian với hệ tọa độ $Oxyz$ cho mặt cầu $(S)\colon x^2+y^2+z^2-4x+2y-2z-3=0$ và điểm $A(5;3;-2)$. Một đường thẳng $d$ thay đổi luôn đi qua $A$ và luôn cắt mặt cầu tại hai điểm phân biệt $M$, $N$. Tính giá trị nhỏ nhất của biểu thức $S=AM+4AN$.
	\loigiai{ 
		\begin{center}
			\begin{tikzpicture}[scale=1,font=\footnotesize,line join=round,line cap=round,>=stealth]
					\def  \r{2}
					 \path
						(0,0) coordinate (I) 
						(160:\r) coordinate  (N)
						(60:\r) coordinate  (M)
						($(M)!(I)!(N)$)  coordinate (H)
						 (180:\r) coordinate (B)
						(intersection of  M--N and I--B) coordinate (A);
					\draw (I) circle (\r); 
					\draw (I)--(A)--(M) 
						 (I)--(H);
					\pic[draw,angle radius=2mm] {right angle = I--H--N}; 
					\foreach \x/\g in {A/100,N/120,M/50,I/-90,H/90}\fill (\x) circle  (1pt)+(\g:.3) node{$\x$};
				\end{tikzpicture}
		\end{center}
		Mặt cầu  $(S)$ có tâm $I(2;-1;1)$, bán kính $R=3$.\\
		$AI=\sqrt{34}>R$, suy  ra $A$ nằm ngoài mặt cầu $(S)$.\\
		Do hai điểm  $M$, $N$ nằm ở vị trí hai đầu một dây cung nên để $S_{\min}$ thì $N$ nằm giữa $A$ và $M$.\\
		Gọi $H$ là trung điểm $MN$, suy ra $IH\perp MN$, $NH=\dfrac{1}2MN$.\\ 
		$S=AH+NH+4(AH-NH)=5AH-3NH$\\ 
		$S=5\sqrt{AI^2-IH^2}-3\sqrt{R^2-IH^2}=5\sqrt{34-x^2}-3\sqrt{9-x^2}$, với $x=IH$.\\ 
		Xét hàm số $f(x)=5\sqrt{34-x^2}-3\sqrt{9-x^2}$ ($0\le  x<3$).\\
		Đạo hàm $f'(x)=\dfrac{-5x}{\sqrt{34-x^2}}+\dfrac{3x}{\sqrt{9-x^2}}=x\left(  \dfrac{-5}{\sqrt{34-x^2}}+\dfrac{3}{\sqrt{9-x^2}} \right)$.\\
		Xét 
		 \allowdisplaybreaks
		\begin{eqnarray*}
			 && \dfrac{-5}{\sqrt{34-x^2}}+\dfrac{3}{\sqrt{9-x^2}}>0\\
			 &\Leftrightarrow& 5\sqrt{9-x^2}<3\sqrt{34-x^2}\\
			&\Leftrightarrow& 225-25x^2<9\cdot 34-9x^2\Leftrightarrow 16x^2+81>0\text{ (luôn đúng)}. 
		\end{eqnarray*}
		Suy  ra $f'(x)\ge 0,\forall x\in [0;3)$, $f'(x)=0\Leftrightarrow x=0\Rightarrow f(x)$ đồng biến trên $[0;3)$.\\
		Suy ra  $\underset{[0;3)}{\mathop{\min }}\,f(x)=f(0)=5\sqrt{34}-9$.
	} 
\end{bt}
